

% new file names from date 20171119: 
% block4_171119.tex 

\renewcommand{\thisfile}{block4\_190328}
 
\cleardoublepage 
\mysection{Dualizing Complexes over Commutative Rings} 
\label{sec:dual-cplx-comm-rng} 

\AYcopyright 

In this section we finally explain what was outlined, as a motivating 
discussion, in Subsection \ref{subsec:2160} of the Introduction. Dualizing 
complexes are perhaps the most compelling reason to study derived categories. In 
the commutative ring setting of the current section the technicalities are 
milder than in the geometric setting (see Remark \ref{rem:4192}) and the 
noncommutative ring setting (that is treated in Sections \ref{sec:BDC} and 
\ref{sec:rigid-DC-NC}). 

In the first and third subsections of this section we talk about {\em dualizing 
complexes} and {\em residue complexes}, respectively, over commutative rings. 
This material is based on the original treatment by A. Grothendieck in 
\cite{RD}, but with a much more detailed discussion. Sandwiched between them is 
a reminder on the  classification of injective modules, due to E. Matlis. 

The last two subsections are on {\em rigid dualizing complexes} in the sense of 
M. Van den Bergh. 

In this section all rings are commutative $\K$-rings by default (see 
Convention \ref{conv:2490}). However, the base ring $\K$ will remain implicit 
most of the time. 


\mysubsection{Dualizing Complexes} \label{subsec:du-cplxs}

Let $A$ be a commutative ring. The category of $A$-modules is 
$\dcat{M}(A) = \cat{Mod} A$. Because $A$ is 
commutative, the Hom bifunctor has its target in $\dcat{M}(A)$~: 
\[ \opn{Hom}_A(-,-) : \dcat{M}(A)^{\mrm{op}} \times \dcat{M}(A)
\to \dcat{M}(A) . \]
Likewise for the right derived bifunctor: 
\[ \opn{RHom}_A(-,-) : \dcat{D}(A)^{\mrm{op}} \times \dcat{D}(A)
\to \dcat{D}(A) . \]
See Proposition \ref{prop:3105}.

Let $M \in \dcat{C}(A)$. The DG $A$-module 
$\opn{End}_A(M) = \opn{Hom}_A(M, M)$ 
is a noncommutative DG central $A$-ring; namely there is a central DG ring 
homomorphism 
\begin{equation} \label{eqn:2390}
\opn{hm}_{M} : A \to \opn{Hom}_A(M, M) 
\end{equation}
called {\em homothety}. 
When we forget the ring structure, $\opn{hm}_{M}$ becomes a homomorphism in
$\dcat{C}_{\mrm{str}}(A)$. 

\begin{dfn} \label{dfn:2392}
Given a complex $M \in \dcat{D}(A)$, the 
{\em derived homothety morphism}%
\index{Homothety morphism! derived}
\[ \opn{hm}^{\mrm{R}}_{M} : A \to \opn{RHom}_A(M, M) \]
is the morphism in $\dcat{D}(A)$ with this formula:
$\opn{hm}^{\mrm{R}}_{M} := \eta^{\mrm{R}}_{M,M} \circ \opn{Q}(\opn{hm}_{M})$.
\end{dfn}

In other words, the diagram 
\[ \UseTips \xymatrix @C=10ex @R=8ex {
A
\ar[r]^(0.35){\opn{Q}(\opn{hm}_{M})}
\ar@(u,u)[rr]^{\opn{hm}^{\mrm{R}}_{M}}
&
\opn{Hom}_A(M, M) 
\ar[r]^(0.45){\eta^{\mrm{R}}_{M,M}}
&
\opn{RHom}_A(M, M) 
} \]
in $\dcat{D}(A)$ is commutative.
The letter ``R'' in the superscript of the expression 
$\opn{hm}^{\mrm{R}}_{M}$ refers to ``right derived functor''. 

\begin{exer} \label{exer:2393}
Prove that if $\rho : M \to I$ is a K-injective resolution, then 
the diagram 
\[ \UseTips \xymatrix @C=10ex @R=6ex {
A
\ar[r]^(0.3){\opn{Q}(\opn{hm}_{I})}
\ar@(dr,l)[drr]_{\opn{hm}^{\mrm{R}}_{M}}
&
\opn{Hom}_A(I, I)
\ar[r]^{\eta^{\mrm{R}}_{I, I}}_{\cong}
&
\opn{RHom}_A(I, I)
\ar[d]^{\opn{RHom}(\opn{Q}(\rho), \opn{Q}(\rho)^{-1})}_{\cong}
\\
&
&
\opn{RHom}_A(M, M)
} \]
in $\dcat{D}(A)$ is commutative.
\end{exer}

\begin{exer} \label{exer:2165}
Formulate and prove a version of the previous exercise with a K-projective 
resolution of $M$. 
\end{exer}

\begin{dfn} \label{dfn:2394}
A complex $M \in \dcat{D}(A)$ is said to have the {\em derived Morita property}
\index{Derived Morita property} 
if the derived homothety morphism 
$\opn{hm}^{\mrm{R}}_{M} : A \to \opn{RHom}_A(M, M)$ 
in $\dcat{D}(A)$ is an isomorphism. 
\end{dfn}

\begin{prop} \label{prop:2335}
The following conditions are equivalent for a complex 
$M \in \dcat{D}(A)$~\tup{:}
\begin{itemize}
\rmitem{i} $M$ has the  derived Morita property.

\rmitem{ii} The homothety ring homomorphism 
$A \to \opn{End}_{\dcat{D}(A)}(M)$
is a bijective, and 
$\opn{Hom}_{\dcat{D}(A)} \bigl( M, M[p] \bigr) = 0$ 
for all $p \neq 0$. 

\rmitem{iii} The $A$-module 
$\opn{H}^0 \bigl( \opn{RHom}_{A}(M, M) \bigr)$ 
is free of rank $1$, with basis the element $\opn{id}_M$, and 
$\opn{H}^p \bigl( \opn{RHom}_{A}(M, M) \bigr) = 0$ 
for every $p \neq 0$.  
\end{itemize}
\end{prop}

\begin{exer}
Prove Proposition \ref{prop:2335}.
(Hint: see Corollary \ref{cor:2120} and the preceding material.) 
\end{exer}

\begin{rem} \label{rem:2390}
In some texts (e.g.\ in \cite{AIL}), a complex $M$ with the derived Morita 
property is called a {\em semi-dualizing complex}. This name is only partly 
justified, because this 
property occurs in the definition of a dualizing complex -- see Definition 
\ref{dfn:2155} below. However, there is a whole other class of 
complexes with the derived Morita property -- these are the {\em tilting 
complexes}. Often these two classes of complexes are disjoint. More on these 
notions, and their noncommutative variants, can be found in Sections 
\ref{sec:perf-tilt-NC} and \ref{sec:rigid-DC-NC} of the book. 
\end{rem}

From here on the commutative ring $A$ is assumed to be noetherian.
Recall that $\dcat{D}_{\mrm{f}}(A)$ is the full subcategory of 
$\dcat{D}(A)$ on the complexes with finitely generated cohomology modules. The 
subcategory $\dcat{D}_{\mrm{f}}(A)$ is triangulated. 

The next definition first appeared in \cite[Section V.2]{RD}. 
The injective dimension of a complex was defined in Definition \ref{dfn:2128}.

\begin{dfn} \label{dfn:2155}
Let $A$ be a noetherian commutative ring. A complex of $A$-modules $R$ is 
called a {\em dualizing complex} 
\index{Dualizing complex! commutative}
if it has the following three 
properties\tup{:} 
\begin{enumerate}
\rmitem{i} $R \in \dcat{D}^{\mrm{b}}_{\mrm{f}}(A)$. 

\rmitem{ii} $R$ has finite injective dimension. 

\rmitem{iii} $R$ has the derived Morita property.
\end{enumerate}
\end{dfn}

Recall that in the traditional literature (e.g.\ \cite{Mats}), a noetherian 
ring $A$ is called {\em regular} if all its local 
rings $A_{\p}$ are regular local rings. The {\em Krull dimension} of $A$ is the 
dimension of the scheme $\opn{Spec}(A)$; namely the supremum of the 
lengths of strictly ascending chains of prime ideals in $A$. 
In practice we never see regular rings that are not finite dimensional (there 
are only pretty exotic examples of them). 
The following definition will simplify matters for us:

\begin{dfn} \label{dfn:3075}
We shall say that a noetherian commutative ring $A$ is 
{\em regular}% 
\index{Regular commutative ring} 
if it has 
finite Krull dimension, and all its local rings $A_{\p}$ are regular local 
rings, in the sense of \cite{Mats}. 
\end{dfn}

Every field $\K$, and the ring of integers $\Z$, are regular rings. 
If $A$ is regular, then so is the polynomial ring 
$A[t_1, \ldots, t_n]$ in $n < \infty$ variables, and also the 
localization of $A$ at every multiplicatively closed set. See 
\cite[Chapter 7]{Mats}. 

As proved by Serre (see \cite[Theorem 19.2]{Mats}) a regular ring $A$ has 
{\em finite global cohomological dimension}. This means that there is a number 
$d \in \N$ such that for all modules $M, N \in \dcat{M}(A)$ and all 
$q > d$, the modules $\opn{Ext}^q_A(M, N)$ vanish. This implies that every
$A$-module $M$ has injective, projective and flat dimensions $\leq d$.
It follows that every complex 
$M \in \dcat{D}^{\mrm{b}}(A)$ has finite injective, projective and flat 
dimensions (see Definitions \ref{dfn:2128} and \ref{dfn:2390}).  

\begin{exa} \label{exa:2171}
Let $A$ be a regular ring. Taking $R := A$ we see that $R$ satisfies condition 
(ii) of Definition \ref{dfn:2155}. The other two conditions hold regardless of 
the regularity of $A$. Thus the complex $R = A$ is a dualizing complex over the 
ring $A$. 

In  Subsection \ref{subsec:2160} of the Introduction we used this fact for 
the ring $A = \Z$. 
\end{exa}

\begin{dfn} \label{dfn:2156}
Given a dualizing complex $R \in \dcat{D}(A)$, the {\em duality functor}
\index{Duality functor} 
associated to it is the triangulated functor 
\[ D : \dcat{D}(A)^{\mrm{op}} \to \dcat{D}(A), \quad 
D := \opn{RHom}_A(-, R) .  \]
\end{dfn}

The notation ``$D$'' deliberately keeps the dualizing complex $R$ implicit. 
Note that upon applying the functor $D$ to the object $A \in \dcat{D}(A)$ we 
get $D(A) = R$. 

Let us choose a K-injective resolution $\rho : R \to I$. There is an 
isomorphism of triangulated functors
\begin{equation} \label{eqn:3235}
\opn{pres}_{I} : D \iso \opn{Hom}_A(-, I)
\end{equation}
from $\dcat{D}(A)^{\mrm{op}}$ to $\dcat{D}(A)$,
that we call a presentation of $D$. For every 
$M \in \dcat{D}(A)$ the diagram 
\begin{equation} \label{eqn:4195}
\UseTips \xymatrix @C=8ex @R=6ex {
D(M)
\ar[r]^(0.4){\opn{id}}
\ar[d]^{\cong}_{\opn{pres}_{I, M}}
&
\opn{RHom}_A(M, R)
\ar[d]^{\opn{RHom}(\opn{id}_M, \opn{Q}(\rho))}_{\cong}
\\
\opn{Hom}_A(M, I)
\ar[r]^{\eta^{\mrm{R}}_{M, I}}_{\cong}
&
\opn{RHom}_A(M, I)
}
\end{equation}
is commutative. (Note that we can choose $I$ to be a bounded complex of 
injectives, by Proposition \ref{prop:2115}.)

Let $M, I \in \dcat{C}(A)$. There is a homomorphism 
\begin{equation} \label{eqn:3236}
\opn{ev}_{M, I} : M \to \opn{Hom}_A \bigl( \opn{Hom}_A(M, I), I \bigr) 
\end{equation}
in $\dcat{C}_{\mrm{str}}(A)$, called {\em Hom-evaluation}, with formula 
\[ \opn{ev}_{M, N}(m)(\phi) := (-1)^{p \cd q} \cd \phi(m) \]
for $m \in M^p$ and $\phi \in \opn{Hom}_A(M, I)^q$. 

\begin{lem} \label{lem:2156} 
Let $R$ be a dualizing complex over $A$, with associated duality functor $D$. 
\begin{enumerate}
\item The functor 
$D \circ D : \dcat{D}(A) \to \dcat{D}(A)$
is triangulated.

\item There is a unique morphism  
$\opn{ev}^{\mrm{R, R}} : \opn{Id} \to D \circ D$ 
of triangulated functors from $\dcat{D}(A)$ to itself, called {\em derived 
Hom-eval\-uation}, such that for every K-injective resolution  
$\rho : R \to I$, and every complex $M \in \dcat{D}(A)$, the diagram 
\[ \UseTips \xymatrix @C=8ex @R=6ex {
M
\ar[r]^(0.4){\opn{ev}^{\mrm{R, R}}_M}
\ar[dr]_{\opn{Q}(\opn{ev}_{M, I}) \ \ }
&
D(D(M))
\ar[d]^{\cong}
\\
&
 \opn{Hom}_A \bigl( \opn{Hom}_A (M, I), I \bigr)
} \]
is commutative. Here the vertical isomorphism is a double application of the 
presentation $\opn{pres}_{I}$.

\item For $M = A$ there is equality 
$\opn{ev}^{\mrm{R, R}}_A = \opn{hm}^{\mrm{R}}_{R}$ 
of morphisms $A \to (D \circ D)(A)$ in $\dcat{D}(A)$.
\end{enumerate}
\end{lem}

\begin{proof}
(1) Choose a K-injective resolution $\rho : R \to I$. 
Let's write 
$F := \lb \opn{Hom}_A \bigl( \opn{Hom}_A (-, I), I \bigr)$.
The functor $F$ can be seen as a DG functor from $\dcat{C}(A)$ to itself, that 
is then made into a functor from $\dcat{D}(A)$ to itself. So by Theorem 
\ref{thm:1265}(1), $F$ is a triangulated functor (i.e.\ there is a 
translation isomorphism $\tau_F$ etc.). Using the presentation 
(\ref{eqn:3235}) twice we get an isomorphism
$D \circ D \cong  F$  of functors from $\dcat{D}(A)$ to itself. This makes 
$D \circ D$ triangulated. 

\medskip \noindent 
(2) With the K-injective resolution $\rho : R \to I$ above, we define 
the morphism $\opn{ev}$ so as to make the diagram commutative. According to
Theorem \ref{thm:1265}(2), $\opn{ev}$ is a morphism of triangulated functors. 
Because the resolution $\rho : R \to I$ is unique up to homotopy, 
we get the same morphism $\opn{ev}$ regardless of choice of resolution. 

\medskip \noindent 
(3) This is clear from Exercise \ref{exer:2393}. 
\end{proof}

\begin{exer} \label{exer:3245}
Suppose that $A$ is a regular ring, and we take the dualizing complex $R := A$. 
Let $P$ be a bounded complex of finitely generated projective $A$-modules. 
Show that:
\begin{enumerate}
\item There is a canonical isomorphism 
$D(P) \cong \opn{Hom}_{A}(P, A)$ 
in $\dcat{D}(A)$. 

\item There is a canonical isomorphism 
\[ (D \circ D)(P) \cong \opn{Hom}_{A} \bigl( \opn{Hom}_{A}(P, A), A \bigr) \]
in $\dcat{D}(A)$.

\item Under the isomorphism in item (2), the morphism 
$\opn{ev}^{\mrm{R, R}}_{P}$ goes to 
$\opn{Q}(\opn{ev}_{P, A})$. 

\item The homomorphism $\opn{ev}_P$ is a quasi-isomorphism. 
\end{enumerate}
\end{exer}

Here is the first important result regarding dualizing complexes. 

\begin{thm} \label{thm:2155}
Suppose $R$ is a dualizing complex over the noetherian commutative ring $A$, 
with associated duality functor $D$. Then for every complex 
$M \in \dcat{D}_{\mrm{f}}(A)$ the following hold\tup{:}
\begin{enumerate}
\item The complex $D(M)$ belongs to $\dcat{D}_{\mrm{f}}(A)$. 

\item The morphism 
\[ \opn{ev}^{\mrm{R, R}}_M : M \to D(D(M)) \]
in $\dcat{D}(A)$ is an isomorphism. 
\end{enumerate}
\end{thm}

\begin{proof}
(1) Condition (ii) of Definition \ref{dfn:2155} says that the functor $D$ has 
finite cohomological dimension. Condition (i) says that 
$D(A) \in \dcat{D}_{\mrm{f}}(A)$.
The assertion follows from Theorem \ref{thm:2160}, with 
$\cat{N}_0 := \dcat{M}_{\mrm{f}}(A)$. 

\medskip \noindent
(2) The composition $D \circ D$ is a triangulated functor with finite 
cohomological dimension (at most twice the injective dimension of $R$). The 
cohomological dimension of the identity functor
$\opn{Id}$ is $0$ (if $A \neq 0$). 
By condition (iii) of Definition \ref{dfn:2155} we know that 
$\opn{ev}_A$ is an isomorphism. 
Now we can use Theorem \ref{thm:2135}. 
\end{proof}

\begin{cor} \label{cor:2155}
Under the assumptions of Theorem \tup{\ref{thm:2155}}, let 
$\star$ be one of the boundedness conditions $\mrm{b}$, $+$, $-$ or 
$\bra{\mrm{empty}}$, and let  $-\star$ be the reversed boundedness condition.
Then the functor 
\[ D : \dcat{D}^{\star}_{\mrm{f}}(A)^{\mrm{op}} \to 
\dcat{D}^{-\star}_{\mrm{f}}(A) \]
is an equivalence of triangulated categories. 
\end{cor}

\begin{proof}
The previous theorem tells us that $D$ is its own quasi-inverse. The claim 
about the boundedness holds because $D$ has finite cohomological dimension.
\end{proof}

We saw that dualizing complexes exist over regular rings. This fact is used 
for the very general existence result Theorem \ref{thm:2170}. But first we need 
some preparation. 

Here are some adjunction properties related to homomorphisms between 
commutative rings. These are enhancements of the material from Subsection 
\ref{subsec:adjs-comm}.

Recall that a ring homomorphism $u : A \to B$ gives rise to a forgetful functor 
$\opn{Rest}_u : \dcat{D}(B) \to \dcat{D}(A)$.
This functor is going to be implicit in the discussion below. 

\begin{lem} \label{lem:2170}
Let $A \to B$ be a ring homomorphism. 
\begin{enumerate}
\item If $I \in \dcat{C}(A)$ is K-injective, then 
$J := \opn{Hom}_A(B, I) \in \dcat{C}(B)$ is K-injective. 

\item Given $M \in \dcat{D}(A)$, let us define 
$N := \opn{RHom}_A(B, M) \in \dcat{D}(B)$.
Then there is an isomorphism 
$\opn{RHom}_B(-, N) \cong \opn{RHom}_A(-, M)$
of triangulated functors
$\dcat{D}(B)^{\mrm{op}} \to \dcat{D}(B)$.
\end{enumerate}
\end{lem}

\begin{proof}
(1) This is an adjunction calculation. Suppose $L \in \dcat{C}(B)$ is acyclic. 
There are isomorphisms  
\begin{equation} \label{eqn:2182}
\opn{Hom}_B(L,J) \cong 
\opn{Hom}_B \bigl( L, \opn{Hom}_A(B, I) \bigr) \cong  
\opn{Hom}_A(L,I) 
\end{equation}
in $\dcat{C}(B)$. Since $I$ is K-injective over $A$, the complex
$\opn{Hom}_A(L,I)$ is acyclic. 

\medskip \noindent
(2) Choose a K-injective resolution $M \to I$ in $\dcat{C}(A)$. 
Let $J := \opn{Hom}_A(B, I)$. Then $N \to J$ is a K-injective 
resolution in 
$\dcat{C}(B)$. There are isomorphisms of triangulated functors 
\begin{equation} \label{eqn:2180}
\opn{RHom}_A(-, M) \cong \opn{Hom}_A(-, I) 
\end{equation}
and 
\begin{equation} \label{eqn:2181}
\opn{RHom}_B(-, N) \cong \opn{Hom}_B(-, J) ,
\end{equation}
where the functors (\ref{eqn:2180}) are contravariant functors from 
$\dcat{D}(A)$ to itself, and the functors (\ref{eqn:2181}) are contravariant 
functors from $\dcat{D}(B)$ to itself. But given $L \in \dcat{C}(B)$
we can view $\opn{Hom}_A(L, I)$ as a complex of $B$-modules, and in this way 
the functors (\ref{eqn:2180}) become contravariant triangulated functors
from $\dcat{D}(B)$ to itself. Formula (\ref{eqn:2182}) shows that the functors 
(\ref{eqn:2180}) and (\ref{eqn:2181}) are isomorphic.
\end{proof}

\begin{lem} \label{lem:2171}
Let $A \to B$ be a flat ring homomorphism, let 
$M \in \dcat{D}^{-}_{\mrm{f}}(A)$, and let 
$N \in \dcat{D}^{+}(A)$. Then there is an isomorphism 
\[ \opn{RHom}_A(M, N) \ot_A B \cong 
\opn{RHom}_B(B \ot_A M, B \ot_A N) \]
in $\dcat{D}(B)$. This isomorphism is functorial in $M$ and $N$. 
\end{lem}

\begin{proof}
Theorem \ref{thm:2700} tells us that 
\[  \opn{RHom}_A(M, N) \ot_A B \cong 
\opn{RHom}_A(M, B \ot_A N) \]
in $\dcat{D}(B)$. Then, by forward adjunction (Theorem \ref{thm:3085})
we get 
\[ \opn{RHom}_A(M, B \ot_A N) \cong \opn{RHom}_B(B \ot_A M, B \ot_A N) . \]
The functoriality is clear. 
\end{proof}

\begin{lem} \label{lem:2185} 
Let $I$ be an $A$-module. The following conditions are equivalent\tup{:}
\begin{enumerate}
\rmitem{i} $I$ is injective.

\rmitem{ii} For every finitely generated $A$-module $M$ the module 
$\opn{Ext}^1_A(M, I)$ is zero.  
\end{enumerate}
\end{lem}

\begin{exer} \label{exer:2185}
Prove Lemma \ref{lem:2185}. (Hint: use the Baer criterion, Theorem 
\ref{thm:151}.) 
\end{exer}

\begin{lem} \label{lem:2172} 
The injective dimension of a complex $N \in \dcat{D}(A)$ equals the 
cohomological dimension of the functor
\[ \opn{RHom}_{A}(-, N)|_{\dcat{M}_{\mrm{f}}(A)^{\mrm{op}}} : 
\dcat{M}_{\mrm{f}}(A)^{\mrm{op}} \to \dcat{D}(A) . \]
\end{lem}

\begin{proof}
By definition the injective dimension of $N$, say $d$,
is the cohomological dimension of the functor
$\opn{RHom}_{A}(-, N) : \dcat{D}(A)^{\mrm{op}} \to \dcat{D}(A)$.
Let $d'$ be  the cohomological dimension of the functor
$\opn{RHom}_{A}(-, N)|_{\dcat{M}_{\mrm{f}}(A)^{\mrm{op}}}$. 
Obviously the inequality $d \geq d'$ holds. 
For the reverse inequality we may assume that $\opn{H}(N)$ is nonzero
and $d' < \infty$. This implies that there are integers $q_1 = q_0 + d'$
such that for every $M \in \dcat{M}_{\mrm{f}}(A)$ there is an inclusion
\[ \opn{con} \bigl( \opn{H} \bigl( \opn{RHom}_{A}(M, N) \bigr) \bigr) \sub 
[q_0, q_1] . \]
In particular, for $M = A$, we get 
$\opn{con}(\opn{H}(N)) \sub [q_0, q_1]$. 
Let $N \to J$ be an injective resolution in 
$\dcat{C}(A)$ with $\opn{inf}(J) \geq q_0$.  
Take $I := \opn{smt}^{\leq q_1}(J)$, the smart truncation from
Definition \ref{dfn:2320}.
The proof of Proposition \ref{prop:2115}, plus Lemma \ref{lem:2185}, 
show that $N \to I$ is an injective resolution. But then 
$\opn{RHom}_{A}(-, N) \cong \opn{Hom}_{A}(-, I)$,
so this functor has cohomological displacement in the interval 
$[q_0, q_1]$, that has length $d'$. 
\end{proof}

Recall that a ring homomorphism $A \to B$ is called {\em finite} if it makes
$B$ into a finitely generated $A$-module.  

\begin{prop} \label{prop:2200}
Let $A \to B$ be a finite ring homomorphism, and let $R_A$ be a dualizing 
complex over $A$. Then the complex 
\[ R_B := \opn{RHom}_{A}(B, R_A) \in \dcat{D}(B) \]
is a dualizing complex over $B$. 
\end{prop}

\begin{proof}
Consider the duality functors
$D_{A} := \opn{RHom}_{A}(-, R_{A})$
and $D_{B} := \lb \opn{RHom}_B(-, R_B)$.
As explained in the proof of Lemma \ref{lem:2170}(2),
they are isomorphic as contravariant triangulated functors from $\dcat{D}(B)$ 
to itself. 
Since $R_{B} = D_A(B)$ and $B \in \dcat{D}^{\mrm{b}}_{\mrm{f}}(A)$,
by Corollary \ref{cor:2155} we have 
$R_B \in \dcat{D}^{\mrm{b}}_{\mrm{f}}(A)$. 
But then also $R_B \in \dcat{D}^{\mrm{b}}_{\mrm{f}}(B)$.
Next, because $D_B(L) \cong D_A(L)$ for all $L \in \dcat{D}(B)$, this implies 
that the cohomological dimension of $D_B$ is at most that of 
$D_A$, which is finite. We see that the injective dimension of the complex 
$R_B$ is finite. Lastly, there is an isomorphism
$D_B \circ D_B \cong D_A \circ D_A$
as functors from $\dcat{D}^{\mrm{b}}_{\mrm{f}}(B)$ to itself, 
and hence 
$\opn{ev}^{\mrm{R, R}} : \opn{Id} \to D_B \circ D_B$ is an isomorphism.
Applying this to the object 
$B \in \dcat{D}^{\mrm{b}}_{\mrm{f}}(B)$
we see that 
$\opn{hm}^{\mrm{R}}_{R_B} = \opn{ev}^{\mrm{R, R}}_{B} : B \to 
(D_B \circ D_{B})(B)$  
is an isomorphism. So $R_B$ has the derived Morita property.
The conclusion is that $R_B$ is a dualizing complex over 
$B$. 
\end{proof}

Recall that a ring homomorphism $A \to B$ is called a {\em localization} if $B$ 
is isomorphic, as an $A$-ring, to the localization $A_S$ of $A$ with respect to 
some multiplicatively closed subset $S$. 

\begin{prop} \label{prop:2201}
Let $A \to B$ be a localization ring homomorphism, and let $R_A$ be a 
dualizing complex over $A$. Then the complex 
\[ R_B := B \ot_A R_A \in \dcat{D}(B) \]
is a dualizing complex over $B$. 
\end{prop}

\begin{proof}
It is clear that 
$R_B \in \dcat{D}^{\mrm{b}}_{\mrm{f}}(B)$.
By Lemma \ref{lem:2172}, to compute the injective dimension of 
$R_B$ it is enough to look at $\opn{RHom}_B(M, R_B)$
for $M \in \dcat{M}_{\mrm{f}}(B)$. 
We can find a finitely generated $A$-submodule $M' \sub M$ such that 
$B \cd M' = M$; and then $M \cong B \ot_{A} M'$. Lemma \ref{lem:2171} 
tells us that 
\[ \opn{RHom}_B(M, R_B) \cong 
\opn{RHom}_{A}(M', R_{A}) \ot_{A} B . \]
We conclude that the injective dimension of $R_B$ is at most that of 
$R_{A}$, which  is finite. Lastly, by the same lemma we get an isomorphism 
\[ \opn{RHom}_B(R_B, R_B) \cong 
\opn{RHom}_{A}(R_{A}, R_{A}) \ot_{A} B  , \]
and it is compatible with the morphisms from $B$. Thus $R_B$ has the derived 
Morita property.  
\end{proof}

Recall that a ring homomorphism $A \to B$ is called {\em finite type} if 
$B$ is finitely generated as an $A$ ring (i.e.\ $B$ is a quotient of the 
polynomial ring $A[t_1, \ldots, t_n]$ for some $n < \infty$). 

\begin{dfn} \label{dfn:3211}
A ring homomorphism $u : A \to B$ is called {\em essentially finite type}
\index{Essentially finite type}
(EFT) if it can be factored into $u = u_{\mrm{loc}} \circ u_{\mrm{ft}}$, where 
$u_{\mrm{ft}} : A \to B_{\mrm{ft}}$ is a finite type ring homomorphism, and 
$u_{\mrm{loc}} : B_{\mrm{ft}} \to B$ is a localization homomorphism.
\end{dfn}

This is illustrated in the following commutative diagram. 
\[ \UseTips \xymatrix @C=14ex @R=6ex {
A
\ar[r]_{\tup{finite type}}^{u_{\mrm{ft}}}
\ar@(u,u)[rr]^{u}
&
B_{\mrm{ft}}
\ar[r]_{\tup{localization}}^{u_{\mrm{loc}}}
&
B 
} \]

\begin{prop} \label{prop:3255}
Let $u : A \to  B$ be an EFT ring homomorphism.
\begin{enumerate}
\item Let $v : B \to C$ be another EFT ring homomorphism. Then 
$v \circ u : A \to C$ is an EFT ring homomorphism.

\item Let $A \to A'$ be a ring homomorphism, and define 
$B' := A' \ot_A B$. Then the induced ring homomorphism 
$u' : A' \to B'$ is EFT. 

\item If the ring $A$ is noetherian, then $B$ is also noetherian. 
\end{enumerate}
\end{prop}

\begin{exer} \label{exer:3256}
Prove Proposition \ref{prop:3255}.
\end{exer}

\begin{exa} \label{exa:3255}
Let $\K$ be a noetherian ring, let $X$ be a finite type $\K$-scheme, and let 
$x \in X$ be a point. Then the local ring $\OO_{X, x}$ is an EFT $\K$-ring. 
\end{exa}

\begin{thm} \label{thm:2170}
Let $\K$ be a regular noetherian commutative ring, and let $A$ be an 
essentially finite type commutative $\K$-ring. Then $A$ has a dualizing 
complex% 
\index{Dualizing complex! commutative}. 
\end{thm}

See Definition \ref{dfn:3075} regarding regular rings. 

\begin{proof}
The ring homomorphism $\K \to A$ can be factored as 
$\K \to A_{\mrm{pl}} \to A_{\mrm{ft}} \to A$, 
where $A_{\mrm{pl}} = \K[t_1, \ldots, t_n]$ is a polynomial ring,
$A_{\mrm{pl}} \to A_{\mrm{ft}}$ is surjective, and $A_{\mrm{ft}} \to A$ is a 
localization. (The subscripts stand for ``polynomial'' and ``finite type'' 
respectively.) According to \cite[Theorem 19.5]{Mats} the ring $A_{\mrm{pl}}$ 
is regular; so, as shown in Example \ref{exa:2171}, the complex
$R_{\mrm{pl}} := A_{\mrm{pl}}$ is a dualizing complex over $A_{\mrm{pl}}$. 

Define 
$R_{\mrm{ft}} := \opn{RHom}_{A_{\mrm{pl}}}(A_{\mrm{ft}}, R_{\mrm{pl}}) \in 
\dcat{D}(A_{\mrm{ft}})$. 
By Proposition \ref{prop:2200} this is a 
dualizing complex over $A_{\mrm{ft}}$. 
Finally define 
$R := A \ot_{A_{\mrm{ft}}} R_{\mrm{ft}} \in \dcat{D}^{\mrm{b}}_{\mrm{f}}(A)$. 
By Proposition \ref{prop:2201} this is dualizing complex over $A$. 
\end{proof}

The proof of Theorem \ref{thm:2170} might give the impression that $A$ could 
have a lot of nonisomorphic dualizing complexes. This is not quite true, as the 
next theorem demonstrates. 

\begin{thm} \label{thm:2175}
Let $A$ be a noetherian commutative ring with connected spectrum, and let $R$ 
and $R'$ be dualizing complexes over $A$. Then there is a rank $1$ projective 
$A$-module $L$ and an integer $d$, such that 
$R' \cong R \ot_A L[d]$ in $\dcat{D}(A)$. 
\end{thm}

We need some lemmas before proving the theorem. 

\begin{lem}[K\"unneth Trick] \label{lem:2177} 
Let $M, M' \in  \dcat{D}^{-}(A)$, and let $i, i' \in \Z$ be such that 
$\opn{sup}(\opn{H}(M)) \leq i$ and $\opn{sup}(\opn{H}(M')) \leq i'$. 
Then 
\[ \opn{H}^{i + i'}(M \ot^{\mrm{L}}_A M') \cong 
\opn{H}^i(M) \ot_A \opn{H}^{i'}(M')  \]
as $A$-modules. 
\end{lem}

\begin{exer} \label{exer:2178}
Prove Lemma \ref{lem:2177}.
\end{exer}

\begin{lem}[Projective Truncation Trick] \label{lem:2189} 
Let $M \in \dcat{D}^-(A)$, with $i_1 := \lb \opn{sup}(\opn{H}(M)) \in \Z$. 
Assume the $A$-module $P := \opn{H}^{i_1}(M)$ is projective. 
Then there is an isomorphism 
$M \cong \opn{smt}^{\leq i_1 - 1}(M) \oplus P[- i_1]$  
in $\dcat{D}(A)$. 
\end{lem}

\begin{exer} \label{exer:2189}
Prove Lemma \ref{lem:2189}. (Hint: first replace $M$ with its truncation 
$\opn{smt}^{\leq i_1}(M)$. Then prove that $P$ is a direct summand of 
$M^{i_1}$.) 
\end{exer}

By a {\em principal open set} in the affine scheme in $\opn{Spec}(A)$ we mean a 
set of the form $\opn{Spec}(A_s)$, where $A_s$ is the localization of $A$ at an 
element $s \in A$. Thus 
\[ \opn{Spec}(A_s) = \{ \p \in \opn{Spec}(A) \mid s \notin \p \} . \]

\begin{lem} \label{lem:2186} 
Let $M, M' \in \dcat{M}_{\mrm{f}}(A)$, and let $\p \sub A$ be a prime ideal. 
\begin{enumerate}
\item If $M_{\p} \neq 0$ and $M'_{\p} \neq 0$ then 
$M_{\p} \ot_{A_{\p}} M'_{\p} \neq 0$. 

\item If $M_{\p} \ot_{A_{\p}} M'_{\p} \cong A_{\p}$ then 
$M_{\p} \cong M'_{\p} \cong A_{\p}$.

\item If $M_{\p} \cong A_{\p}$, then there is a principal open neighborhood 
$\opn{Spec}(A_s)$ of $\p$ in $\opn{Spec}(A)$ such that 
$M_s \cong A_s$ as $A_s$-modules. 
\end{enumerate}
\end{lem}

\begin{exer} \label{exer:2187}
Prove Lemma \ref{lem:2186}. (Hint: use the Nakayama Lemma.) 
\end{exer}

Here is a pretty difficult technical lemma. 

\begin{lem} \label{lem:2175}
In the situation of the theorem, let 
$M, M' \in \dcat{D}^{-}_{\mrm{f}}(A)$
satisfy \lb 
$M \ot^{\mrm{L}}_A M' \cong A$ in $\dcat{D}(A)$. 
Then $M \cong L[d]$ in $\dcat{D}(A)$ for some rank $1$ projective $A$-module 
$L$ and an integer $d$.
\end{lem}

\begin{proof}
For each prime $\p \sub A$ let $M_{\p} := A_{\p} \ot_A M$, and define 
$e_{\p} := \opn{sup}(\opn{H}(M_{\p}))$ in $\Z \cup \{ -\infty \}$.
Define the generalized integer $e'_{\p}$ similarly. 

Fix one prime $\p$. 
The associativity and symmetry of the left derived tensor product imply the 
existence of these isomorphisms
\begin{equation} \label{2189}
\begin{aligned}
& M_{\p} \ot^{\mrm{L}}_{A_{\p}} M'_{\p} = 
(A_{\p} \ot^{\mrm{L}}_{A} M) \ot^{\mrm{L}}_{A_{\p}}
(A_{\p} \ot^{\mrm{L}}_{A} M')
\\
& \qquad 
\cong A_{\p} \ot^{\mrm{L}}_{A} (M \ot^{\mrm{L}}_A M') 
\cong A_{\p} \ot^{\mrm{L}}_{A} A \cong A_{\p} 
\end{aligned}
\end{equation}
in $\dcat{D}(A_{\p})$. Since $A_{\p} \neq 0$, it follows that
$\opn{H}(M_{\p}) \neq 0$ and $\opn{H}(M'_{\p}) \neq 0$.
So $e_{\p}, e'_{\p} \in \Z$, and 
$\opn{H}^{e_{\p}}(M_{\p})$, $\opn{H}^{e'_{\p}}(M'_{\p})$ are nonzero finitely 
generated $A_{\p}$-modules. By Lemma \ref{lem:2186}(1) we know that 
\[ \opn{H}^{e_{\p}}(M_{\p}) \ot_{A_{\p}} \opn{H}^{e'_{\p}}(M'_{\p}) \neq 0 . \]
According to Lemma \ref{lem:2177} we have 
\[ \opn{H}^{e_{\p}}(M_{\p}) \ot_{A_{\p}} \opn{H}^{e'_{\p}}(M'_{\p}) \cong 
\opn{H}^{(e_{\p} + e'_{\p})}(M_{\p} \ot^{\mrm{L}}_{A_{\p}} M'_{\p}) \cong 
\opn{H}^{(e_{\p} + e'_{\p})}(A_{\p}) . \]
But $A_{\p}$ is concentrated in degree $0$; this forces 
$e_{\p} + e'_{\p} = 0$ and 
\[ \opn{H}^{e_{\p}}(M_{\p}) \ot_{A_{\p}} \opn{H}^{e'_{\p}}(M'_{\p}) \cong 
A_{\p} \]
in $\dcat{D}(A_{\p})$. 
By Lemma \ref{lem:2186}(2) we now see that
$\opn{H}^{e_{\p}}(M_{\p}) \cong \opn{H}^{e'_{\p}}(M'_{\p}) \cong A_{\p}$. 
According to Lemma \ref{lem:2189} there are isomorphisms 
\begin{equation} \label{eqn:2196}
M_{\p} \cong A_{\p}[-e_{\p}] \oplus \opn{smt}^{\leq e_{\p} - 1}(M_{\p}) 
\end{equation}
and 
\[ M'_{\p} \cong A_{\p}[-e'_{\p}] \oplus \opn{smt}^{\leq e'_{\p} - 1}(M'_{\p})  
\]
in $\dcat{D}(A_{\p})$. 
These, with (\ref{2189}), give an isomorphism
\begin{equation} \label{eqn:2195}
\bigl( A_{\p}[-e_{\p}] \oplus \opn{smt}^{\leq e_{\p} - 1}(M_{\p}) \bigr) 
\ot^{\mrm{L}}_{A_{\p}}
\bigl( A_{\p}[-e'_{\p}] \oplus \opn{smt}^{\leq e'_{\p} - 1}(M'_{\p}) \bigr)
\cong A_{\p} .
\end{equation}
The left side of (\ref{eqn:2195}) is the direct sum of four objects. 
Passing to the cohomology of (\ref{eqn:2195}) we see that 
$N := \opn{H} \bigl( \opn{smt}^{\leq e_{\p} - 1}(M_{\p})[-e'_{\p}] \bigr)$ 
is a direct summand of $A_{\p}$.
But, since $e'_{\p} + e_{\p} = 0$, the graded module $N$ is concentrated in the 
degree interval $[\infty,  -1]$. It follows that $N = 0$. Therefore, by 
(\ref{eqn:2196}), we deduce that 
\begin{equation} \label{eqn:2197}
M_{\p} \cong A_{\p}[-e_{\p}] . 
\end{equation}

The calculation above works for every prime $\p$. 
From (\ref{eqn:2197}) we get 
\begin{equation} \label{eqn:2189}
A_{\p} \ot_A \opn{H}^i(M) \cong \opn{H}^i(M_{\p}) \cong 
\begin{cases}
A_{\p} & \tup{if} \ i = e_{\p} , 
\\
0 & \tup{otherwise} . 
\end{cases}
\end{equation}

We now use Lemma \ref{lem:2186}(3) to deduce that for every prime $\p$ there is 
an open neighborhood $U_{\p}$ of $\p$ in $\opn{Spec}(A)$ such that 
$\opn{H}^{e_{\p}}(M_{\q}) \cong A_{\q}$ for all $\q \in U_{\p}$.
This implies, by equation (\ref{eqn:2189}), that 
$e_{\q} = e_{\p}$.  Therefore $\p \mapsto e_{\p}$ is a locally constant function
$\opn{Spec}(A) \to \Z$. We assumed that $\opn{Spec}(A)$ is connected, and this 
implies that this is a constant function, say $e_{\p} = -d$ for some integer 
$d$. 

Define $L := \opn{H}^{-d}(M) \in \dcat{M}_{\mrm{f}}(A)$. 
Using truncation we see that $M \cong L[d]$ in 
$\dcat{D}(A)$. We know that $L_{\p} \cong A_{\p}$ for all primes $\p$. 
Finally, Lemma \ref{lem:2171} says that the $A$-module $L$ is projective. 
\end{proof}

\begin{rem} \label{rem:2185}
Lemma \ref{lem:2175} is actually true in much greater generality: 
the ring $A$ does not have to be noetherian, and we do not have to assume 
that the complexes $M$ and $M'$ have bounded above or finite cohomology. 
The proof is harder (see the proof of Theorem \ref{thm:3585}).
\end{rem}

\begin{proof}[Proof of Theorem \tup{\ref{thm:2175}}]
Define the duality functors
$D := \opn{RHom}_{A}(-, R)$ and \lb $D' := \opn{RHom}_{A}(-, R')$; these are 
finite 
dimensional contravariant triangulated functors from $\dcat{D}_{\mrm{f}}(A)$ to 
itself. And define $F := D' \circ D$ and $F' := D \circ D'$, that are finite 
dimensional (covariant) triangulated functors from $\dcat{D}_{\mrm{f}}(A)$ to 
itself. Let
\begin{equation} \label{eqn:2177}
M := F(A) = D'(D(A)) = \opn{RHom}_{A}(R, R')
\end{equation}
and
\[ M' := F'(A) = D(D'(A)) = \opn{RHom}_{A}(R', R) . \]
These are objects of $\dcat{D}^{\mrm{b}}_{\mrm{f}}(A)$. 

For every object $N \in \dcat{D}(A)$ there is a morphism 
\[ \psi_N : N \ot^{\mrm{L}}_A \opn{RHom}_{A}(R, R') \to 
\opn{RHom}_{A} \bigl( \opn{RHom}_{A} (N, R) , R' \bigr) \]
defined as follows: we choose a K-projective resolution $P \to N$
and a K-injective resolution $R' \to I'$. Then $\psi_N$ is represented by the 
obvious homomorphism of complexes 
\[ P \ot^{}_A \opn{Hom}_{A}(R, I') \to 
\opn{Hom}_{A} \bigl( \opn{Hom}_{A} (P, R) , I' \bigr) . \]
As $N$ changes, $\psi_N$ is a morphism of triangulated functors
$\psi : (-) \ot^{\mrm{L}}_A M \to D' \circ D  = F$.
For $N = A$ the morphism $\psi_A$ is an isomorphism, by equation 
(\ref{eqn:2177}). The functor $F$ has finite cohomological dimension, and the 
functor $(-) \ot^{\mrm{L}}_A M$ has bounded above cohomological displacement. 
According to Theorem \ref{thm:2135}, the morphism 
$\psi_N$ is an isomorphism for every $N \in \dcat{D}^-_{\mrm{f}}(A)$. 
In particular this is true for $N := M'$. So, using Theorem \ref{thm:2155}, 
we obtain 
\[ M' \ot^{\mrm{L}}_A M \cong (D' \circ D)(M') \cong 
(D' \circ D \circ D \circ D')(A) \cong A . \]
According to Lemma \ref{lem:2175} there is an isomorphism 
$M \cong L[d]$. Finally, using the isomorphism $\psi_R$, we get 
\[ R \ot_A L[d] \cong F(R) = D'(D(R)) \cong D'(A) = R' . \qedhere \]
\end{proof}

What if $\opn{Spec}(A)$ has more than one connected component?

\begin{dfn} \label{dfn:3210}
Let $A$ be a ring. 
\begin{enumerate}
\item We say that $A$ has a {\em finite connected component decomposition}
if 
\[ \opn{Spec}(A) = \coprod_{i = 1}^{m} \, \opn{Spec}(A_i) , \] 
a finite disjoint union of nonempty connected closed subschemes. 

\item If $A$ has a finite connected component decomposition as above, then 
the 
{\em connected component decomposition}%
\index{Connected component decomposition}
of $A$ is the ring isomorphism 
$A \cong \prod_{i = 1}^{m} \, A_i$.
\end{enumerate}
\end{dfn}

Of course, if $A$ has a finite connected component decomposition, then 
this decomposition is unique, up to renumbering. 

\begin{prop} \label{prop:5080}
If $A$ is noetherian, then it has a finite connected component decomposition.
\end{prop}

\begin{exer} \label{exer:5080} \mbox{}
\begin{enumerate}
\item Prove the proposition above. 

\item Find a ring $A$ that does not have a finite connected component 
decomposition.
\end{enumerate}
\end{exer}

\begin{cor} \label{cor:2175}
Let $R$ and $R'$ be dualizing complexes over $A$, and let 
$A = \prod_{i = 1}^m A_i$ be the connected component decomposition of $A$.
Then there is an isomorphism 
\[ R' \cong R \ot_A \bigl( L_1[d_1] \oplus \cdots \oplus L_m[d_m] \bigr) \]
in $\dcat{D}(A)$, where each $L_i$ is a rank $1$ projective 
$A_i$-module, and each $d_i$ is an integer. Furthermore, the modules $L_i$ are 
unique up to isomorphism, and the integers $d_i$ are unique. 
\end{cor}

\begin{exer} \label{exer:2176}
Prove Corollary \ref{cor:2175}.
\end{exer}

\begin{rem} \label{rem:2195}
A rank $1$ projective $A$-module $L$ is also called an {\em invertible 
$A$-module}. This is because $L$ is invertible for the tensor product. 
Recall that the group of isomorphism classes of invertible $A$-modules is 
the {\em commutative Picard group} $\opn{Pic}_A(A)$. 

The {\em commutative derived Picard group} $\opn{DPic}_A(A)$ is the abelian 
group $\opn{Pic}_A(A) \times \Z^m$ that classifies dualizing complexes over 
$A$, as in Corollary \ref{cor:2175}.

Now assume that $A$ is {\em noncommutative}, and flat central over a 
commutative ring $\K$. There are noncommutative versions of dualizing complexes 
and of ``invertible'' complexes, that are called {\em tilting complexes}. The 
latter form the nonabelian group $\opn{DPic}_{\K}(A)$, and it classifies 
noncommutative dualizing complexes. See
\cite{Ric1}, \cite{Ric2}, \cite{Kel1}, \cite{Ye4} and \cite{RoZi}. 
We shall study this material in Sections \ref{sec:perf-tilt-NC} 
and \ref{sec:rigid-DC-NC} of the book. 
\end{rem}

\begin{rem} \label{rem:2196}
The lack of uniqueness of dualizing complexes has always been a source of 
difficulty. A certain uniqueness or functoriality is needed, already for 
proving existence of dualizing complexes on schemes. 

In \cite{RD} Grothendieck utilized local and global duality in 
order to formulate a suitable uniqueness of dualizing 
complexes. This approach was very cumbersome (even without providing details!).

Since then there have been a few approaches in the literature to attack this 
difficulty. Generally speaking, these approaches came in two flavors:
\begin{itemize}
\item {\em Representability}. This started with P. Deligne's Appendix to 
\cite{RD}, and continued most notably in the work of A. Neeman, J. Lipman and 
their coauthors. See \cite{Ne2}, \cite{Li2} and their references. 

\item {\em Explicit Constructions}. Mostly in the early work of Lipman et 
al., including \cite{Li1} and \cite{LNS}, and in the work of A. Yekutieli 
\cite{Ye2}, and \cite{Ye3} and \cite{Ye5}. 
\end{itemize}

In Subsection \ref{subsec:rig-res-cplx-rng} of the book we will present {\em 
rigid residue complexes}, for which there is a built-in uniqueness, and even 
functoriality (see Remark \ref{rem:4191}). 
\end{rem}

\newpage



\mysubsection{Interlude: The Matlis Classification of Injective Modules}
\label{subsec:mr-inj-res}

We start with a few facts about injective modules over rings that are neither 
commutative nor noetherian. Sources for this material are \cite{Rot} and
\cite{Lam}. 

\begin{dfn} \label{dfn:2205} \mbox{}
\begin{enumerate}
\item Let $M$ be an $A$-module. A submodule $N \sub M$ is called an 
{\em essential submodule} 
\index{Essential! submodule}
if for every nonzero submodule $L \sub M$, the 
intersection $N \cap L$ is nonzero. In this case we also say that $M$ is an 
{\em essential extension} of $N$. 

\item An {\em essential monomorphism}
\index{Essential! monomorphism} 
is a monomorphism $\phi : N \inj M$
whose image is an essential submodule of $M$.   

\item Let $M$ be an $A$-module. An {\em injective hull}
\index{Injective hull}
(or {\em injective 
envelope}) of $M$ is an injective module $I$, together with an essential 
monomorphism $M \inj I$. 
\end{enumerate}
\end{dfn}

\begin{prop} \label{prop:2225}
Every $A$-module $M$ admits an injective hull. 
\end{prop}

\begin{proof}
See \cite[Theorem 3.30]{Rot} or \cite[Section 3.D]{Lam}. 
\end{proof}

There is a weak uniqueness result for injective hulls.

\begin{prop} \label{prop:2207}
Let $M$ be an $A$-module, and suppose $\phi : M \inj I$ and 
$\phi' : M \inj I'$ are monomorphisms into injective modules. 
\begin{enumerate}
\item If $\phi$ is essential, then there is a monomorphism 
$\psi : I \inj I'$ such that $\psi \circ \phi = \phi'$.

\item If $\phi'$ is also essential, then $\psi$ above is an isomorphism.
\end{enumerate}
\end{prop}

\begin{exer} \label{exer:2207}
Prove Proposition \ref{prop:2207}.
\end{exer}

In classical homological algebra we talk about the minimal injective resolution 
of a module $M$. Let us recall it. We start by taking the injective hull
$\phi : M \inj I^0$. This gives an exact sequence 
\[ 0 \to M \xar{\phi} I^0 \to M^1 \to 0 , \]
where $M^1$ is the cokernel of $\phi$. Then we take the injective hull
$M^1 \inj I^1$, and this gives a longer exact sequence 
\[ 0 \to M \to I^0 \to I^1 \to M^2 \to 0 , \]
and so on. We want to generalize this idea to complexes. 

\begin{dfn} \label{dfn:2210}
\mbox{}
\begin{enumerate}
\item A {\em minimal complex of injective $A$-modules}%
\index{Complex in abelian category! minimal injective}
is a bounded below complex of injective modules $I$, such that for every integer 
$q$ the submodule of cocycles $\opn{Z}^q(I) \sub I^q$ is essential. 

\item Let $M \in \dcat{D}^+(A)$. A 
{\em minimal injective resolution}% 
\index{Resolution! minimal injective}
of $M$ is a quasi-iso\-morphism $M \to I$ into a minimal complex of injectives 
$I$.
\end{enumerate}
\end{dfn}

\begin{prop} \label{prop:2210}
Let $M \in \dcat{D}^+(A)$.
\begin{enumerate}
\item  There exists a minimal injective resolution $\phi : M \to I$.

\item If $\phi' : M \to I'$ is another minimal injective resolution, then there 
is an isomorphism $\psi : I \to I'$ in $\dcat{C}_{\mrm{str}}(A)$ such that 
$\phi' = \psi \circ \phi$.

\item If $M$ has finite injective dimension, then it has a bounded minimal 
injective resolution $I$. 
\end{enumerate}
\end{prop}

\begin{proof}
(1) We know that there is a quasi-isomorphism $M \to J$ where $J$ is a bounded 
below complex of injective modules. For every $q$ let $E^q$ be an injective 
hull of $\opn{Z}^q(J)$. By Proposition \ref{prop:2207}(1) we can assume that 
$E^q$ sits inside $J^q$ like this: $\opn{Z}^q(J) \sub E^q \sub J^q$.
Since $E^q$ is injective, we can decompose $J^q$ into a direct sum:
$J^q \cong E^q \oplus K^q$. The homomorphism 
$\d_J^q : K^q \to J^{q + 1}$ 
is a monomorphism since $K^q \cap \opn{Z}^q(J) = 0$.
And the image $\d_J^q(K^q)$ is contained in $E^{q + 1}$. 
Thus $\d_J^q(K^q)$ is a direct summand of $E^{q + 1}$, and this shows that the 
quotient
\[ I ^{q + 1} := E^{q + 1} / \d_J^q(K^q) \cong
J^{q + 1} / (K^{q + 1} \oplus \d_J^q(K^q)) \]
is an injective module. The canonical surjection of graded modules 
$\pi : J \to I$ is a homomorphism of complexes, with kernel the acyclic complex 
\[ \bigoplus_{q \in \Z} \ \Bigl( K^q[-q] \xar{\d_J^q} \d_J^q(K^q)[-q - 1] 
\Bigr) . \]
Therefore $\pi$ is a quasi-isomorphism. A short calculation shows that $I$ is a 
minimal complex of injectives, i.e.\ $\opn{Z}^q(I) \sub I^q$ is essential. 

\medskip \noindent
(2) See next exercise. (We will not need this fact.) 

\medskip \noindent
(3) According to Proposition \ref{prop:2115}, the complex $J$ that appears in 
item (1) can be chosen to be bounded. 
\end{proof}

\begin{exer} \label{exer:2215}
Prove Proposition \ref{prop:2210}(2). 
\end{exer}

\begin{rem} \label{rem:2210}
Important: the isomorphisms $\psi$ in Propositions \ref{prop:2207} and 
\ref{prop:2210} are not unique (see next exercise). We will see below (in 
Subsection \ref{subsec:rig-res-cplx-rng}) 
that a rigid residue complex is a minimal complex of injectives that has no 
nontrivial rigid automorphisms. 
\end{rem}

\begin{exer} \label{exer:2210}
Take $A := \K[[t]]$, the power series ring over a field $\K$. 
Let $M := A / (t)$, the trivial module (the residue field viewed as an 
$A$-module). 
\begin{enumerate}
\item Find a minimal injective resolution $M \to  I$. 

\item Find a nontrivial automorphism of the complex $I$ in 
$\dcat{C}_{\mrm{str}}(A)$ that fixes the submodule $M \sub I^0$. 
\end{enumerate}
\end{exer}

Now we add the noetherian condition. 

\begin{prop} \label{prop:2205}
Assume $A$ is a left noetherian ring. 
Let $\{ I_z \}_{z \in Z}$ be a collection of injective $A$-modules. Then 
$I := \bigoplus_{z \in Z} I_z$ is an injective $A$-module. 
\end{prop}

\begin{exer} \label{exer:2205}
Prove Proposition \ref{prop:2205}. (Hint: use the Baer criterion.)
\end{exer}

From here on in this subsection all rings are noetherian commutative. For 
them much more can be said.

Recall that a module $M$ is called {\em indecomposable} if it is not the 
direct sum of two nonzero modules.  

\begin{dfn} \label{dfn:2208}
Let $\a \sub A$ be an ideal.
\begin{enumerate}
\item Let $M$ be an $A$-module. The {\em 
$\a$-torsion submodule} of $M$ is the submodule $\Ga_{\a}(M)$ 
consisting of the elements that are annihilated by powers of $\a$. Thus 
\[ \Ga_{\a}(M) = \lim_{i \to} \, \opn{Hom}_A(A / \a^i, M) \sub M . \]

\item If $\Ga_{\a}(M) = M$ then $M$ is called an {\em $\a$-torsion module}. 

\item The functor 
$\Ga_{\a} : \dcat{M}(A) \to \dcat{M}(A)$
is called the {\em $\a$-torsion functor}. 
\end{enumerate}
\end{dfn}

Here are some important properties of the torsion functor.

\begin{prop} \label{prop:3085}
Let $\a$ be an ideal in $A$. 
\begin{enumerate}
\item The functor $\Ga_{\a}$ is left exact. 

\item The functor $\Ga_{\a}$ commutes with infinite direct sums. 
\end{enumerate}
\end{prop}

\begin{exer} \label{exer:3085}
Prove the proposition above. 
\end{exer}

Perhaps the most important theorem about injective modules over 
noetherian commutative rings is the following structural result due to E. 
Matlis \cite{Matl} from 1958. See also \cite[Section V.4]{Ste}, 
\cite[Sections 3.F and 3.I]{Lam}, \cite[Section 18]{Mats} and \cite{BrSh}.

For a prime ideal $\p \sub A$ we write 
$\kk(\p) := A_{\p} / \p_{\p}$, the residue field of $\p$. 

\begin{thm}[Matlis] \label{thm:2026}
Let $A$ be a noetherian commutative ring. 
\begin{enumerate}
\item Let $\p \sub A$ be a prime ideal, and let $J(\p)$ be the injective hull 
of the $A_{\p}$-module $\kk(\p)$. Then, as an $A$-module, $J(\p)$ is
injective, indecomposable and $\p$-torsion. 

\item Suppose $I$ is an indecomposable injective $A$-module. Then 
$I \cong J(\p)$ for a unique prime ideal $\p \sub A$. 

\item Every injective $A$ module $I$ is a direct sum of indecomposable 
injective $A$-modules. 
\end{enumerate}
\end{thm}

Theorem \ref{thm:2026} tells us that every injective $A$-module $I$ can be 
written as a direct sum 
\begin{equation} \label{eqn:2215}
I \cong \bigoplus_{\p \in \opn{Spec}(A)} \, J(\p)^{\oplus \mu_{\p}} 
\end{equation}
for a collection of cardinal numbers 
$\{ \mu_{\p} \}_{\p \in \opn{Spec}(A)}$, called the {\em Bass numbers}. 
General counting tricks can show that the multiplicity $\mu_{\p}$ is an 
invariant of $I$. But we can be more precise: 

\begin{prop} \label{prop:2208}
Let $I$ be an injective $A$-module, with direct sum decomposition 
\tup{(\ref{eqn:2215})}.
Then for every $\p$ there is equality
\[ \mu_{\p} = \opn{rank}_{\kk(\p)} \bigl( 
\opn{Hom}_{A_{\p}} \bigl( \kk(\p) , A_{\p} \ot_A I \bigr) \bigr) . \]
\end{prop}

\begin{proof}
Consider another prime $\q$. If $\q \nsubseteq \p$ then there is an element
$a \in \q - \p$, and then $a$ is both invertible and locally nilpotent on 
$A_{\p} \ot_A J(\q)$. This implies that $A_{\p} \ot_A J(\q) = 0$.
On the other hand, if $\q \sub \p$, then 
$A_{\p} \ot_A J(\q) \cong J(\q)$. Therefore 
$A_{\p} \ot_A I \cong  \bigoplus_{\q \sub \p} J(\p)^{\oplus \mu_{\p}}$.

Next, if $\q \varsubsetneq \p$, then there is an element $b \in \p - \q$, 
and it is both invertible and zero on the module 
$\opn{Hom}_{A_{\p}} \bigl( \kk(\p) , J(\q)  \bigr)$.
The implication is that this module is zero.

Finally, if $\q = \p$ then we have 
\[ \opn{Hom}_{A_{\p}} \bigl( \kk(\p) , J(\p)  \bigr) \cong 
\opn{Hom}_{A_{\p}} \bigl( \kk(\p) , \kk(\p)  \bigr) \cong 
\kk(\p) , \]
because the inclusion $\kk(\p) \sub J(\p)$ is essential. 

Since $\kk(\p)$ is a finitely generated $A_{\p}$-module, the functor 
$\opn{Hom}_{A_{\p}}(\kk(\p), -)$
commutes with infinite direct sums. Therefore, putting all these cases 
together, we see that 
\[ \opn{Hom}_{A_{\p}} \bigl( \kk(\p) , A_{\p} \ot_A I \bigr) \bigr) \cong 
\kk(\p)^{\oplus \mu_{\p}}  \]
as $\kk(\p)$-modules.
\end{proof}




\mysubsection{Residue Complexes} \label{subsec:res-cplx-ring}

In this subsection $A$ is a noetherian commutative ring. 
Here we introduce residue complexes (called residual complexes in \cite{RD}). 
Most of the material is taken from the original \cite{RD}. 
In Example \ref{exa:2245} we will see the relation between the geometry of 
$\opn{Spec}(A)$ and the structure of dualizing complexes over $A$ (continuing 
Example \ref{exa:2252} from the Introduction). 

\begin{lem} \label{lem:2200}
Let $R$ be a dualizing complex over $A$ and let $\p \sub A$ be a prime ideal. 
We write 
$R_{\p} := A_{\p} \ot_A R \in \dcat{D}(A_{\p})$.
There is an integer $d$ such that 
\[ \opn{Ext}^i_{A_{\p}} \bigl(\kk(\p), R_{\p} \bigr) \cong 
\begin{cases}
\kk(\p) & \tup{if} \ i = -d , 
\\
0 & \tup{otherwise} . 
\end{cases} \]
\end{lem}

\begin{proof}
By Proposition \ref{prop:2201}, $R_{\p}$ is a dualizing complex over the local 
ring $A_{\p}$. And by Proposition \ref{prop:2200},
$S := \opn{RHom}_{A_{\p}} \bigl( \kk(\p), R_{\p} \bigr)$ 
is a dualizing complex over the residue field $\kk(\p)$.  
Since $\kk(\p)$ is a field, it is a regular ring, and so it is a dualizing 
complex over itself. Theorem  \ref{thm:2175} tells us that 
$S \cong \kk(\p)[d]$ in $\dcat{D}(\kk(\p))$ 
for some integer $d$.  
\end{proof}

\begin{dfn} \label{dfn:2200}
The number $d$ in Lemma \ref{lem:2200} is called the {\em dimension of $\p$ 
relative to $R$}, and is denoted by $\opn{dim}_R(\p)$. 
In this way we obtain a function
\[  \opn{dim}_R : \opn{Spec}(A) \to  \Z , \]
called the {\em dimension function} 
\index{Dimension function}
associated to $R$. 
\end{dfn}

Let us recall a few notions regarding the combinatorics of prime ideals in a 
ring $A$. A prime ideal $\q$ is an {\em immediate specialization} of another 
prime $\p$ if $\p \subsetneqq \q$, and there is no other prime $\p'$ such that 
$\p \subsetneqq \p' \subsetneqq \q$. 
In other words, if the dimension of the local ring $A_{\q} / \p_{\q}$ is $1$. 

A {\em chain of prime ideals} in $A$ is a sequence 
$(\p_0, \ldots, \p_n)$ of primes such that 
$\p_i \subsetneqq \p_{i + 1}$ for all $i$. The number $n$ is the {\em length} 
of the chain. The chain is called {\em saturated} if for each $i$ the prime 
$\p_{i + 1}$ is an immediate specialization of $\p_i$. 

\begin{thm} \label{thm:2205}
Let $R$ be a dualizing complex over $A$ and let $\p, \q \sub A$ be prime 
ideals. Assume that $\q$ is an immediate specialization of $\p$. Then
\[ \opn{dim}_R(\q) = \opn{dim}_R(\p) - 1 . \]
\end{thm}

To prove this theorem we need a baby version of local cohomology: 
codimension $1$ only. 

Let $\a$ be an ideal in $A$. 
The torsion functor $\Ga_{\a}$ has a right derived functor
$\mrm{R} \Ga_{\a}$. For every complex $M \in \dcat{D}(A)$, the module 
$\opn{H}^p_{\a}(M) := \opn{H}^p(\mrm{R} \Ga_{\a}(M))$ 
is called the {\em $p$-th cohomology of $M$ with support in $\a$}. 
In case $A$ is a local ring and $\m$ is its maximal ideal, then 
$\opn{H}^p_{\m}(M)$ is also called the {\em local cohomology of $M$}. 

Now suppose $\a$ is a principal ideal in $A$, generated by an 
element $a$. Let $A_a  = A[a^{-1}]$ be the localized ring.
For any $A$-module $M$ we write $M_a = A_a \ot_A M$. 
There is a canonical exact sequence 
\begin{equation} \label{eqn:2206}
0 \to \Ga_{\a}(M) \to M \to M_a . 
\end{equation}

\begin{lem} \label{lem:2220}
Let $\a = (a)$ be a principal ideal in $A$. 
\begin{enumerate}
\item For every injective module $I$ the sequence 
$0 \to \Ga_{\a}(I) \to I \to I_a \to 0$
is exact. 

\item For every $M \in \dcat{D}^+(A)$ there is a long exact sequence of 
$A$-modules 
\[ \cdots \to \opn{H}^p_{\a}(M) \to \opn{H}^p(M) \to \opn{H}^p(M_a) \to 
\opn{H}^{p + 1}_{\a}(M) \to  \cdots . \]
\end{enumerate}
\end{lem}

\begin{proof}
(1) Let $J(\q)$ be an indecomposable injective $A$-module. According to Theorem 
\ref{thm:2026}(1), if $a \in \q$ then 
$\Ga_{\a}(J(\q)) = J(\q)$ and $J(\q)_a = 0$. But if $a \notin \q$
then $J(\q) = J(\q)_a$ and $\Ga_{\a}(J(\q)) = 0$.
By Theorem \ref{thm:2026} and Proposition \ref{prop:3085}(2)
we see that each injective module $I$ breaks up into 
a direct sum $I = \Ga_{\a}(I) \oplus I_a$, and this proves that the sequence 
is split exact.  

\medskip \noindent
(2) Choose a resolution $M \to I$ by a bounded below complex of injectives. 
We obtain an exact sequence of complexes as in item (1). The long 
exact sequence in cohomology
\[ \cdots \to \opn{H}^p(\Ga_{\a}(I)) \to \opn{H}^p(I) \to 
\opn{H}^p(I_a) \to \opn{H}^{p + 1}(\Ga_{\a}(I)) \to  \cdots  \]
is what we want. 
\end{proof}

\begin{lem} \label{lem:2235}
Suppose $A$ is an integral domain, with fraction field $K$, such that
$A \neq K$. Then $K$ is not a finitely generated $A$-module. 
\end{lem}

\begin{proof}
Let $a \in A$ be a nonzero element that is not invertible. Then 
\[ A \subsetneqq a^{-1} \cd A \subsetneqq a^{-2} \cd A \subsetneqq
\cdots \sub K \]
is an infinite ascending sequence of $A$-submodules of $K$. 
\end{proof}

\begin{lem} \label{lem:2221}
For every ideal $\a$ and every $M \in \dcat{D}(A)$ there is an isomorphism of 
$A$-modules 
$\opn{H}^p_{\a}(M) \cong \lim_{k \to} \opn{Ext}^p_A( A / \a^k, M)$.
\end{lem}

\begin{proof}
Choose a K-injective resolution $M \to I$. Then, using the fact that cohomology 
commutes with direct limits, we have
\[ \begin{aligned}
& \opn{H}^p_{\a}(M) \cong \opn{H}^p(\Ga_{\a}(I)) \cong 
\opn{H}^p \bigl( \lim\nolimits_{k \to} \, \opn{Hom}_A(A / \a^k, I) \bigr)
\\ & \quad \cong  
\lim_{k \to} \, \opn{H}^p \bigl( \opn{Hom}_A (A / \a^k, I) \bigr) 
\cong \lim_{k \to} \, \opn{Ext}^p_A(A / \a^k, M) .
\end{aligned} \]
\end{proof}

\begin{lem} \label{lem:2250}
Assume $A$ is local, with maximal ideal $\m$. Let $R$ be a dualizing complex 
over $A$, and let $d := \opn{dim}_R(\m)$. Then for every $i \neq -d$ 
the local cohomology $\opn{H}^{i}_{\m}(R)$ vanishes. 
\end{lem}

See Remark \ref{rem:2250} for more about $\opn{H}^{-d}_{\m}(R)$. 

\begin{proof}
We know that 
\[ \opn{Ext}^i_{A} \bigl( \kk(\m), R \bigr) \cong 
\begin{cases}
\kk(\m) & \tup{if} \ i = -d , 
\\
0 & \tup{otherwise} .
\end{cases} \]
Let $N$ be a finite length $A$-module. Since $N$ is gotten from the residue 
field $\kk(\m)$ by finitely many extensions, induction on the length of $N$ 
shows that 
$\opn{Ext}^i_{A}(N, R) = 0$ for all $i \neq -d$. 
This holds in particular for $N := A / \m^k$. 
Now use Lemma \ref{lem:2221}.
\end{proof}


\begin{proof}[Proof of Theorem \tup{\ref{thm:2205}}]
Define $d := \opn{dim}_R(\q)$ and $e := \opn{dim}_R(\p)$.
We need to prove that $e = d + 1$. 

By definition, $d$ is the unique integer s.t.\ 
$\opn{Ext}^{-d}_{A_{\q}} \bigl( \kk(\q), R_{\q} \bigr) \neq 0$.
Let's define $\bar{A} := A / \p$. 
We know that 
$\bar{R} := \opn{RHom}_{A}(\bar{A}, R)$
is a dualizing complex over $\bar{A}$. There are isomorphisms 
\[ \begin{aligned}
& \opn{RHom}_{A_{\q}} \bigl( \kk(\q), R_{\q} \bigr) \cong 
\opn{RHom}_{\bar{A}_{\q}} \bigl( \kk(\q), 
\opn{RHom}_{A_{\q}} \bigl( \bar{A}_{\q} , R_{\q} \bigr) \bigr)
\\
& \quad 
\cong \opn{RHom}_{\bar{A}_{\q}} \bigl( \kk(\q), \bar{R}_{\q} \bigr)
\end{aligned} \]
in $\dcat{D}(\bar{A}_{\q})$, coming from adjunction for the homomorphism 
$A_{\q} \to \bar{A}_{\q}$. 
There is also an isomorphism 
\[ \opn{RHom}_{A_{\p}} \bigl( \kk(\p), R_{\p} \bigr)  
\cong \opn{RHom}_{\bar{A}_{\p}} \bigl( \kk(\p), \bar{R}_{\p} \bigr) \]
in $\dcat{D}(\bar{A}_{\p})$.
Hence we can replace $A$ and $R$ with $\bar{A}_{\q}$ and $\bar{R}_{\q}$ 
respectively. 

Now we have $\p = 0$ and $A = A_{\q}$. Thus $A$ is a $1$-dimensional local 
integral domain, 
with only two primes ideals: $0 = \p$ and the maximal ideal $\q$. Take 
any nonzero element $a \in \q$. Then the localization $A_a$ is the field of 
fractions of $A$, i.e.\ $A_a = \kk(\p)$. 
On the other hand, letting $\a := (a) \sub A$, the quotient 
$A / \a$ is an artinian local ring. So $A / \a$ is a finite length 
$A$-module, the ideal $\a$ is $\q$-primary, and $\Ga_{\a} = \Ga_{\q}$. 
 
By Lemma \ref{lem:2220} there is an exact sequence of $A$-modules 
\[ \cdots \to \opn{H}^{-e}_{\a}(R) \to \opn{H}^{-e}(R) \xar{\phi} 
\opn{H}^{-e}(R_a) \to \opn{H}^{-e + 1}_{\a}(R) \to  \cdots . \]
Since $a \neq 0$ there are equalities $A_a = A_{\p} = \opn{Frac}(A) = \kk(\p)$. 
Then 
$\opn{H}^{-e}(R_a) \cong \kk(\p)$, and this is not a finitely generated 
$A$-module by Lemma \ref{lem:2235}. On the other hand the $A$-module 
$\opn{H}^{-e}(R)$ is finitely generated. 
We conclude that homomorphism $\phi$ is not surjective, and thus 
$\opn{H}^{-e + 1}_{\a}(R) \neq 0$. 
But 
$\opn{H}^{-e + 1}_{\a}(R) = \opn{H}^{-e + 1}_{\q}(R)$,
so according to Lemma \ref{lem:2250} we must have 
$-e + 1 = -d$. Thus $e = d + 1$ as claimed. 
\end{proof}

\begin{cor} \label{cor:2235}
If $A$ has a dualizing complex, then the Krull dimension of $A$ is finite. 
More precisely, if $R$ is a dualizing complex over $A$, then 
$\opn{dim}(A)$ is at most the injective dimension of $R$. 
\end{cor}

\begin{proof}
Let $[i_0, i_1]$ be the injective concentration of the complex $R$. See 
Definition \ref{dfn:2128}. This is a bounded interval. Since 
\[ \opn{Ext}^i_{A_{\p}}(\kk(\p), R_{\p}) \cong
\opn{Ext}^i_{A}(A / \p, R)_{\p} ,  \]
we see that 
$\opn{dim}_{R}(\p) \in -[i_0, i_1] = [-i_1, -i_0]$. 

Let $(\p_0, \ldots, \p_n)$ be a chain of prime ideals in $A$. Because $A$ is 
noetherian, we can squeeze more primes into this chain, until after finitely 
many steps it becomes saturated. According to Theorem \ref{thm:2205}, for a 
saturated chain we have 
$n = \opn{dim}_{R}(\p_0) - \opn{dim}_{R}(\p_n)$.
Therefore $n \leq i_1 - i_0$. 
\end{proof}

\begin{dfn} \label{dfn:2237}
The ring $A$ is called {\em catenary} if for every pair of primes 
$\p \sub \q$ there is a number $n_{\p, \q}$ such that for every saturated chain 
$(\p_0, \ldots, \p_n)$ with $\p_0 = \p$ and $\p_n = \q$, there is equality
$n = n_{\p, \q}$. 
\end{dfn}

\begin{cor} \label{cor:2236}
If $A$ has a dualizing complex, then it is catenary.
\end{cor}

\begin{proof}
Let $R$ be a dualizing complex over $A$. The proof of the previous corollary 
shows that the number 
$n_{\p, \q}= \opn{dim}_{R}(\p) - \opn{dim}_{R}(\q)$
has the desired property. 
\end{proof}

\begin{exa} \label{exa:2245}
This is a continuation of Example \ref{exa:2252} from the Introduction. 
Consider the ring 
$A = \mbb{R}[t_1, t_2, t_3] / (t_3 \cd t_1, \, t_3 \cd t_2)$.
The affine algebraic variety
$X = \opn{Spec}(A) \sub \mbf{A}^3_{\mbb{R}}$
is shown in figure \ref{fig:2250}. 
It is the union of a plane $Y$ and a line 
$Z$, meeting at the origin. 

\begin{figure}
\centering
\includegraphics[scale=0.3]{fig2.jpg}
\caption{An algebraic variety $X$ that is connected but not equidimensional: it 
has irreducible components $Y$ and $Z$ of dimensions $2$ and $1$ respectively.
The generic points $\p \in Y$ and $\q \in Z$ are shown.}
\label{fig:2250}
\end{figure}

Since the ring $A$ is finite type over the field $\R$, it has a dualizing 
complex $R$. We will now prove that there is some integer $i$ s.t.\ 
$\mrm{H}^i(R)$ and $\mrm{H}^{i+1}(R)$ are nonzero. 

Define the prime ideals 
$\m := (t_1, t_2, t_3)$, 
$\q := (t_1, t_2)$ and 
$\p := (t_3)$. Thus $\m$ is the origin, $\q$ is the generic point of the line
$Z = \opn{Spec}(A / \q)$, and $\p$ is the generic point of the plane
$Y = \opn{Spec}(A / \p)$. 
By translating $R$ as needed, we can assume that 
$\opn{dim}_R(\m) = 0$. Since $\m$ is an immediate specialization of $\q$, 
Theorem \ref{thm:2205} tells us that $\opn{dim}_R(\q) = 1$. 
Similarly, since every line in $Y$ passing through the origin gives 
rise to a saturated chain 
$(\p, \q', \m)$, we see that $\opn{dim}_R(\p) = 2$.

Since $\q$ is the generic point of $Z$, its local ring is the residue field: 
$A_{\q} = \kk(\q)$. We know that $\opn{dim}_R(\q) = 1$. Hence 
\[ \kk(\q) \cong \opn{Ext}^{-1}_{A_{\q}}(\kk(\q), R_{\q}) =
\opn{Ext}^{-1}_{A_{\q}}(A_{\q}, R_{\q}) \cong 
\opn{Ext}^{-1}_{A}(A, R)_{\q} \cong \opn{H}^{-1}(R)_{\q} . \]
Therefore $\opn{H}^{-1}(R) \neq 0$. 
A similar calculation involving $\p$ shows that 
$\opn{H}^{-2}(R) \neq 0$.
\end{exa}

\begin{exa} \label{2255}
Let $A$ be a local ring, with maximal ideal $\m$ and residue field  
$\kk(\m)$. Recall that $A$ is called {\em Gorenstein} if the free 
module $A$ has finite injective dimension. The ring $A$ is called called {\em 
Cohen-Macaulay} if its depth is equal to its dimension, where the depth of $A$ 
is the minimal integer $i$ such that 
$\opn{Ext}^i_A(\kk(\m), A) \neq 0$. It is known that Gorenstein implies 
Cohen-Macaulay. See \cite{Mats} for details. 

As is our usual practice (cf.\ Definition \ref{dfn:3075}) we shall say that 
a noetherian commutative ring $A$ is Cohen-Macaulay (resp.\ Gorenstein) if it 
has finite Krull dimension, and all its local rings $A_{\p}$ are Cohen-Macaulay 
(resp.\ Gorenstein) local rings, as defined above. 

Assume $A$ has a connected spectrum, and let $R$ be a dualizing complex over 
$A$. Grothendieck showed in \cite[Section V.9]{RD} that $A$ is a Cohen-Macaulay 
ring iff $R \cong L[d]$ for some finitely generated module $L$ and some integer 
$d$; the proof is not easy. It is however pretty easy to prove (using Theorem 
\ref{thm:2175}) that $A$ is a Gorenstein ring iff $R \cong L[d]$ for some 
invertible module $L$ and some integer $d$. 

There is a lot more to say about the relation between the CM (Cohen-Macaulay) 
property and duality. See Remark \ref{rem:2270}
\end{exa}

Recall that for each $\p \in \opn{Spec}(A)$ we denote by $J(\p)$ the 
corresponding indecomposable injective $A$-module. 

\begin{dfn} \label{dfn:2236}
A {\em residue complex} 
\index{Residue complex}
over $A$ is a complex of $A$-modules $\KK$ having these properties:
\begin{itemize}
\rmitem{i} $\KK$ is a dualizing complex. 

\rmitem{ii} For every integer $d$ there is an isomorphism of $A$-modules 
\[ \KK^{-d} \cong 
\bigoplus_{\substack{ \p \in \opn{Spec}(A) \\  
\opn{dim}_{\KK}(\p)  = d }}  J(\p) \ . \]
\end{itemize}
\end{dfn}

The reason we like residue complexes is this:

\begin{thm} \label{thm:2240}
Suppose $\KK$ and $\KK'$ are residue complexes over $A$ that have the same 
dimension function. Then the homomorphism 
\[ \opn{Q} : \opn{Hom}_{\dcat{C}_{\mrm{str}}(A)}(\KK, \KK') \to 
\opn{Hom}_{\dcat{D}(A)}(\KK, \KK') \]
is bijective. 
\end{thm}

In other words, for each morphism $\psi : \KK \to \KK'$ in 
$\dcat{D}(A)$ there is a unique strict homomorphism of complexes 
$\phi : \KK \to \KK'$ such that $\psi = \opn{Q}(\phi)$. 

\begin{proof}
Since the complex $\KK'$ is K-injective, by Theorem \ref{thm:3135}
we know that the homomorphism 
\[ \opn{Q} : \opn{Hom}_{\dcat{K}(A)}(\KK, \KK') \to 
\opn{Hom}_{\dcat{D}(A)}(\KK, \KK') \]
is bijective. And by definition the homomorphism 
\[ \opn{P} : \opn{Hom}_{\dcat{C}_{\mrm{str}}(A)}(\KK, \KK') \to 
\opn{Hom}_{\dcat{K}(A)}(\KK, \KK') \]
is surjective. It remains to prove that 
$\opn{Hom}_{A}(\KK, \KK')^{-1} = 0$,
i.e.\ here are no nonzero degree $-1$ homomorphisms $\ga : \KK \to \KK'$. 

The residue complexes $\KK$ and $\KK'$ decompose into indecomposable summands 
by the formula in property (ii) of Definition \ref{dfn:2236}. 
A homomorphism $\ga : \KK \to \KK'$ of degree $-1$ 
is nonzero iff at least one of its components 
$\ga_{\p, \q} : J(\p) \to J(\q)$ 
is nonzero, for some $J(\p) \sub \KK^{-i}$ and $J(\q) \sub \KK'^{\, -i - 1}$.
Denoting by $\opn{dim}$ the dimension function of both these dualizing 
complexes, we have $\opn{dim}(\p) = i$ and $\opn{dim}(\q) = i + 1$. 
But the lemma below says that $\q$ has to be a specialization of $\p$. 
Therefore, as in the proof of Corollary \ref{cor:2235}, there is an inequality 
in the oppose direction:  
$\opn{dim}(\p) \geq \opn{dim}(\q)$. 
We see that it is impossible to have a  nonzero degree $-1$ 
homomorphism $\ga : \KK \to \KK'$. 
\end{proof}

\begin{lem} \label{lem:2240}
Let $\p, \q$ be prime ideals. If there is a nonzero homomorphism 
$\ga : J(\p) \to J(\q)$, then 
$\q$ is a specialization of $\p$. 
\end{lem}

\begin{proof}
Assume $\q$ is not a specialization of $\p$; i.e.\ 
$\p \nsubseteq \q$. So there is an element $a \in \p - \q$. 
Let $\ga : J(\p) \to J(\q)$ be a homomorphism, and consider the module 
$N := \ga(J(\p)) \sub J(\q)$. Since $J(\p)$ is $\p$-torsion, the element $a$ 
acts on $N$ locally-nilpotently. On the other hand, $J(\q)$ is a module over 
$A_{\q}$, so $a$ acts invertibly on $J(\q)$, and hence it has zero annihilator 
in $N$. The conclusion is that $N = 0$. 
\end{proof}

Here is a general existence theorem. 
Minimal injective resolutions were defined in Definition \ref{dfn:2210}.
Their existence was proved in Proposition \ref{prop:2210}. 

\begin{thm} \label{thm:2245}
Suppose the ring $A$ has a dualizing complex $R$. 
Let $R \to \KK$ be a minimal injective resolution of $R$. Then $\KK$ is a 
residue complex over $A$. 
\end{thm}

The proof is after two lemmas. 

\begin{lem} \label{lem:2230}
Let $S \sub A$ be a multiplicatively closed set, with localization $A_S$. 
For each complex of $A$-modules $M$ we write $M_S := A_S \ot_A M$.
\begin{enumerate}
\item If $I$ is an injective $A$-module, then $I_S$ is an 
injective $A_S$-module.

\item If $I$ is an injective $A$-module and $M \sub I$ is an essential 
$A$-submodule, then $M_S \sub I_S$ is an essential $A_S$-submodule.

\item If $I$ is a minimal complex of injective $A$-modules, then 
$I_S$ is a minimal complex of injective $A_S$-modules.
\end{enumerate}
\end{lem}

\begin{proof}
(1) By Theorem \ref{thm:2026} there is a direct sum decomposition 
$I \cong I' \oplus I''$, where 
$I' \cong \bigoplus_{\p \cap S = \varnothing}  
J(\p)^{\oplus \mu_{\p}}$
and 
$I'' \cong \bigoplus_{\p \cap S \neq \varnothing} 
J(\p)^{\oplus \mu_{\p}}$.
If $\p \cap S = \varnothing$ then $J(\p) \cong J(\p)_S$ is an injective 
$A_S$-module; and if $\p \cap S \neq \varnothing$ then $J(\p)_S = 0$. 
We see that $I_S \cong I'$ is  an injective $A_S$-module.

\medskip \noindent 
(2) Denote by $\la : I \to I_S$ the canonical homomorphism. 
Under the decomposition $I \cong I' \oplus I''$ above, 
$\la|_{I'} : I' \to I_S$ is an isomorphism. 

Let $L$ be a nonzero $A_S$-submodule of $I_S$. 
Consider the $A$-submodule 
$L' := (\la|_{I'})^{-1}(L) \sub I'$;
so $\la|_{L'} : L' \to L$ is an isomorphism and $L' \neq 0$.  
Because $M \sub I$ is essential, the intersection $M \cap L'$ is 
nonzero. But then $\la(M \cap L')$ is a nonzero submodule of 
$M_S \cap L$, so $M_S \cap L$ is nonzero. 

\medskip \noindent 
(3) By part (1) the complex $I_S$ is a bounded below complex of injective 
$A_S$-modules. Exactness of localization shows that 
$\opn{Z}^n(I_S) = \opn{Z}^n(I)_S$ inside $I^n_S$; so by part (2) the inclusion 
$\opn{Z}^n(I_S) \inj I^n_S$ is essential.
\end{proof}

\begin{lem} \label{lem:2231}
Let $\a \sub A$ be an ideal, and define $B := A / \a$. 
\begin{enumerate}
\item If $I$ is an injective $A$-module, then 
$J := \opn{Hom}_A(B, I)$ is an injective $B$-module. 

\item Let $I$ and $J$ be as above. If $M \sub I$ is an essential $A$-submodule, 
then $N := \opn{Hom}_A(B, M)$ is an essential $B$-submodule of $J$.

\item If $I$ is a minimal complex of injective $A$-modules, then 
$J := \opn{Hom}_A(B, I)$ is a minimal complex of injective $B$-modules,
\end{enumerate}
\end{lem}

\begin{proof}
(1) This is immediate from adjunction. 

\medskip \noindent
(2) We identify $J$ and $N$ with the submodules of $I$ and $M$ respectively 
that are annihilated by $\a$. Let $L \sub J$ be a nonzero $B$-submodule.
Then $L$ is a nonzero $A$-submodule of $I$. 
Because $M$ is essential, the intersection $L \cap M$ is nonzero.
But $L \cap M$ is annihilated by $\a$, so it sits inside $N$, and in fact 
$L \cap M = L \cap N$. 

\medskip \noindent
(3) By part (1) the complex $J$ is a bounded below complex of injective 
$B$-modules. Left exactness of $\opn{Hom}_A(B, -)$ shows that 
$\opn{Z}^n(J) = \opn{Hom}_A(B, \opn{Z}^n(I))$ inside $J^n$; so by part (2) 
the inclusion $\opn{Z}^n(J) \inj J^n$ is essential.
\end{proof}

\begin{proof}[Proof of Theorem \tup{\ref{thm:2245}}]
Since $\KK \cong R$ in $\dcat{D}(A)$ it follows that $\KK$ is a dualizing 
complex. To show that $\KK$ has property (ii) of Definition \ref{dfn:2236} we 
have to count multiplicities. For every $\p$ and $d$ let $\mu_{\p, d}$ be the 
multiplicity of $J(\p)$ in $\KK^{-d}$, so that  
\[ \KK^{-d} \cong 
\bigoplus_{\p \in \opn{Spec}(A)}  J(\p)^{\oplus \mu_{\p, d}} . \]
We have to prove that 
\begin{equation} \label{eqn:2245}
\mu_{\p, d} = 
\begin{cases}
1 & \tup{if} \ \opn{dim}_{\KK}(\p) = d ,  
\\
0 & \tup{otherwise} . 
\end{cases}
\end{equation}
Now by Lemma \ref{lem:2230}(3) the complex $\KK_{\p} = A_{\p} \ot_A \KK$ is a  
minimal complex of injective  $A_{\p}$-modules. Because $\KK_{\p}$ is 
K-injective over $A_{\p}$ we get   
\[ \opn{Ext}^{-d}_{A_{\p}}(\kk(\p), R_{\p}) \cong
\opn{H}^{-d} \bigl( \opn{Hom}_{A_{\p}}(\kk(\p), \KK_{\p}) \bigr)  \]
as $\kk(\p)$-modules. By Lemma \ref{lem:2231}(3) the complex 
$\opn{Hom}_{A_{\p}}(\kk(\p), \KK_{\p})$
is a  minimal complex of injective  $\kk(\p)$-modules.
It is easy to see (and we leave this verification to the reader) that a 
minimal complex of injectives over a field must have zero differential. 
Therefore 
\[ \opn{H}^{-d} \bigl( \opn{Hom}_{A_{\p}}(\kk(\p), \KK_{\p}) \bigr) \cong 
\opn{Hom}_{A_{\p}}(\kk(\p), \KK_{\p}^{-d}) . \]
Now by arguments like in the proofs of Lemmas \ref{lem:2230}(1) 
and \ref{lem:2240} we know that 
\[ \opn{Hom}_{A_{\p}} \bigl( \kk(\p), J(\q)_{\p} \bigr) \cong 
\begin{cases}
\kk(\p) & \tup{if} \ \q = \p ,  
\\
0 & \tup{otherwise} . 
\end{cases} \]
It follows that 
\[ \opn{Hom}_{A_{\p}} \bigl( \kk(\p), \KK_{\p}^{-d} \bigr) \cong 
\kk(\p)^{\oplus \mu_{\p, d}} . \]
We see that 
\[ \opn{rank}_{\kk(\p)} \bigl( \opn{Ext}^{-d}_{A_{\p}}(\kk(\p), R_{\p}) \bigr) 
= \mu_{\p, d} . \]
But by Definition \ref{dfn:2200} this number satisfies (\ref{eqn:2245}). 
\end{proof}

\begin{cor} \label{cor:2245}
If $\KK$ is a residue complex over $A$, then it is a minimal complex 
of injective $A$-modules. 
\end{cor}

\begin{proof}
Let $\phi : \KK \to \KK'$ be a minimal injective resolution of $\KK$. According 
to Theorem \ref{thm:2245}, $\KK'$ is also a residue complex. Now 
$\opn{Q}(\phi) : \KK \to \KK'$ is an isomorphism in $\dcat{D}(A)$, so by Theorem
\ref{thm:2240} we know that $\phi : \KK \to \KK'$ is an isomorphism in 
$\dcat{C}_{\mrm{str}}(A)$. 
\end{proof}

\begin{exer} \label{exer:2245}
Find a direct proof of Corollary \ref{cor:2245}, without resorting to Theorems 
\ref{thm:2245} and \ref{thm:2240}. (Hint: look at the proof of Proposition 
\ref{prop:2210}.) 
\end{exer}

We end this subsection with three remarks. 

\begin{rem} \label{rem:2245}
Here is a brief explanation of {\em Matlis Duality}. For more details see 
\cite[Section V.5]{RD}, \cite[Theorem 18.6]{Mats} or \cite[Section 10.2]{BrSh}. 
Assume $A$ is a complete local ring with maximal ideal $\m$. As usual, the 
category of finitely generated $A$-modules is $\dcat{M}_{\mrm{f}}(A)$. There is 
also the category $\dcat{M}_{\mrm{a}}(A)$ of artinian $A$-modules. These are 
full abelian subcategories of $\dcat{M}(A)$. Note that these subcategories are 
characterized by dual properties: the objects of $\dcat{M}_{\mrm{f}}(A)$ are 
noetherian, i.e.\ they satisfy the ascending chain condition; and the objects 
of $\dcat{M}_{\mrm{a}}(A)$ satisfy the descending chain condition. 

Consider the indecomposable injective module $J(\m)$.  
The functor 
$D := \lb \opn{Hom}_A \bigl( -, J(\m) \bigr)$  
is exact of course. Matlis Duality asserts that
$D : \dcat{M}_{\mrm{f}}(A)^{\mrm{op}} \to \dcat{M}_{\mrm{a}}(A)$ 
is an equivalence, with quasi-inverse 
$D : \dcat{M}_{\mrm{a}}(A)^{\mrm{op}} \to \dcat{M}_{\mrm{f}}(A)$.

Later in this book we present a noncommutative graded version of Matlis 
Duality -- this is Theorem \ref{thm:4045}.
\end{rem}

\begin{rem} \label{rem:2250}
We now provide a brief discussion of {\em Local Duality}, based on 
\cite[Section V.6]{RD}. (There is a weaker variant of this result, 
for modules instead of complexes, that can be found in 
\cite[Theorem 11.2.6]{BrSh}.) Again $A$ is local, with maximal ideal $\m$. 
Let $R$ be a dualizing complex over $A$. By translating $R$ we can assume that 
$\opn{dim}_R(\m) = 0$. Lemma \ref{lem:2250} tells us that 
$\opn{H}^i_{\m}(R) = 0$ for all $i \neq 0$. A calculation, that relies on 
Matlis duality, shows that 
$\opn{H}^0_{\m}(R) \cong J(\m)$, the indecomposable injective corresponding to 
$\m$. 

Let us fix an isomorphism $\be : \opn{H}^0_{\m}(R) \iso J(\m)$. This induces 
a morphism 
\begin{equation} \label{eqn:2260}
\th_M : \mrm{R} \Ga_{\m}(M) \to 
\opn{Hom}_A \bigl( \opn{RHom}_A(M, R), J(\m) \bigr) ,
\end{equation}
functorial in $M \in \dcat{D}^+(A)$. The Local Duality Theorem
\cite[Theorem V.6.2]{RD} says that 
$\th_M$ is an isomorphism if $M \in \dcat{D}^+_{\mrm{f}}(A)$.

Here is a modern take on this theorem: we can construct the morphism $\th_M$ 
for all $M \in \dcat{D}(A)$. Let's replace $R$ by the residue complex $\KK$ 
(the minimal injective resolution of $R$). Then $\be$ 
is just an $A$-module isomorphism $\be : \KK^0 \iso J(\m)$. 
For each complex $M$ we choose a K-injective resolution $M \to I(M)$.
Then $\th_M$ is represented by the homomorphism 
\[ \til{\th}_M : \Ga_{\m}(I(M)) \to 
\opn{Hom}_A \bigl( \opn{Hom}_A \bigl( I(M), \KK \bigr), \KK^0 \bigr) \]
in $\dcat{C}_{\mrm{str}}(A)$
that sends an element $u \in \Ga_{\m} \bigl( I(M) \bigr)^p$ 
and a homomorphism 
$\phi \in \opn{Hom}_A \bigl( I(M), \KK \bigr)^{-p}$
to $(-1)^p \cd \phi(u) \in \KK^0$. 

We know that the functor $\mrm{R} \Ga_{\m}$ has finite cohomological 
dimension, \lb bounded by the number of generators of the ideal $\m$; see 
\cite{LC} or \cite{PSY}. The functor 
$\opn{RHom}_{A}(-, R)$
has finite cohomological dimension, which is the injective dimension of $R$. 
And the functor 
$\opn{Hom}_{A}(-, J(\m))$ has cohomological dimension $0$. 
Since $A \in \dcat{D}^+_{\mrm{f}}(A)$, the local duality theorem from \cite{RD}
tells us that $\th_A$ is an isomorphism. Now we can apply Theorem \ref{thm:2135}
to conclude that $\th_M$ is an isomorphism for every 
$M \in \dcat{D}_{\mrm{f}}(A)$.

Finally, let us mention that in Subsection \ref{subsec:bal-NCDC-LocDu} there is 
a noncommutative graded version of Local Duality. 
\end{rem}

\begin{rem} \label{rem:2270}
Here is more on the CM (Cohen-Macaulay) property and duality. 
Let $A$ be a noetherian ring with connected spectrum. Assume $A$ has a 
dualizing complex $R$, and corresponding dimension function $\opn{dim}_{R}$. 

Consider a complex $M \in  \dcat{D}_{\mrm{f}}^{\mrm{b}}(A)$.
In \cite{RD} Grothendieck defines $M$ to be a {\em CM complex with respect to 
$R$} if for every prime ideal $\p \sub A$ and every 
$i \neq - \opn{dim}_{R}(\p)$ the local cohomology satisfies 
$\opn{H}^i_{\p}(M_{\p}) = 0$.

It is proved in \cite{RD} that when $A$ is a regular ring, 
$R = A$, and $M$ is a finitely generated $A$-module, then $M$ is a CM module 
(in the conventional sense, see \cite{Mats}) iff it is a CM complex. 

Let $\dcat{D}_{\mrm{f}}^{0}(A)$ be the full subcategory of 
$\dcat{D}_{\mrm{f}}^{\mrm{b}}(A)$ on the complexes $M$ such that 
$\opn{H}^i(M) = 0$ for all $i \neq 0$.
We know that $\dcat{D}_{\mrm{f}}^{0}(A)$
is equivalent to 
$\dcat{M}_{\mrm{f}}(A) = \cat{Mod}_{\mrm{f}} A$.
In \cite{YeZh2} it was proved that the following
are equivalent for a complex $M \in  \dcat{D}_{\mrm{f}}^{\mrm{b}}(A)$~:
\begin{enumerate}
\rmitem{i} The complex $M$ is CM w.r.t.\ $R$. 

\rmitem{ii} The complex $\opn{RHom}_A(M, R)$ belongs to 
$\dcat{D}_{\mrm{f}}^{0}(A)$. 
\end{enumerate}

It follows that the CM complexes form an abelian subcategory of 
$\dcat{D}^{\mrm{b}}_{\mrm{f}}(A)$,
dual to $\dcat{M}_{\mrm{f}}(A)$. 
In fact, they are the heart of a perverse t-structure on 
$\dcat{D}^{\mrm{b}}_{\mrm{f}}(A)$,
and hence they deserve to be called {\em perverse finitely generated 
$A$-modules}. Geometrically, on the scheme $X := \opn{Spec}(A)$, 
the CM complexes inside $\dcat{D}^{\mrm{b}}_{\mrm{c}}(X)$
form a stack of abelian categories, and so they are 
{\em perverse coherent sheaves}. All this is explained in 
\cite[Section 6]{YeZh2}.
\end{rem}



\mysubsection{Van den Bergh Rigidity}
\label{subsec:squaring}

As we saw in Theorem \ref{thm:2175}, a dualizing complex $R$ 
over a noetherian commutative ring $A$ is not unique. This lack of uniqueness 
(not to mention any sort of functoriality!) was the source of major 
difficulties 
in \cite{RD}, first for gluing dualizing complexes on schemes, and then for 
producing the trace morphisms associated to maps of schemes. 

In 1997, M. Van den Bergh \cite{VdB} discovered the idea of {\em rigidity} for 
dualizing complexes. This was done in the context of noncommutative ring 
theory: $A$ is a noncommutative noetherian ring, central over a base field 
$\K$. The theory of noncommutative rigid dualizing complexes was developed 
further in several papers of J.J. Zhang and A. Yekutieli, among them 
\cite{YeZh1} and \cite{YeZh2}. We shall talk about this noncommutative theory in 
Section \ref{sec:rigid-DC-NC} of the book. 

Here we will deal with the commutative side only, which turns out to be 
extremely powerful. Before explaining it, let us first observe that this is 
one of the rare cases in which an idea originating from noncommutative 
algebra had significant impact in commutative algebra and algebraic geometry.  

In this subsection we study Van den Bergh rigidity in the following context: 

\begin{setup} \label{set:2285}
$A$ is a commutative ring, and $B$ is a flat commutative $A$-ring. 
\end{setup}

We introduce the notion of {\em rigid complex over $B$ relative to $A$},
and describe some of its properties. In Subsection 
\ref{subsec:rig-res-cplx-rng} we will discuss {\em rigid dualizing and residue
complexes}. This material is adapted from the papers \cite{YeZh3}, \cite{YeZh4} 
and \cite{Ye11}. 

The theory of rigid complexes does not really require the assumption
that $B$ is flat over $A$, but flatness makes the theory much easier. See 
Remark \ref{rem:2285}. 

Consider the enveloping ring $B \ot_A B$. It comes equipped with a few ring 
homomorphisms: 
\begin{equation} \label{eqn:2281}
B \xar{\eta_i} B \ot_{A} B \xar{\ep} B ,
\end{equation}
where $\eta_0(b) := b \ot 1$, $\eta_1(b) := 1 \ot b$, and  
$\ep(b_0 \ot b_1) := b_0 \cd b_1$. We view $B$ as a module over 
$B \ot_{A} B$ via $\ep$. Of course $\ep \circ \eta_i = \opn{id}_B$. 

Suppose we are given $B$-modules $M_0$ and $M_1$. Then the tensor product
$M_0 \ot_{A} M_1$ is a $(B \ot_{A} B)$-module. In this way we get an additive 
bifunctor
\[ (- \ot_A -) : \dcat{M}(B) \times \dcat{M}(B) \to \dcat{M}(B \ot_{A} B) . \]
Passing to complexes, and then to homotopy categories, we obtain a triangulated 
bifunctor
\begin{equation} \label{eqn:2276}
(- \ot_A -) : \dcat{K}(B) \times \dcat{K}(B) \to \dcat{K}(B \ot_{A} B) . 
\end{equation}

\begin{lem} \label{lem:2276}
The bifunctor \tup{(\ref{eqn:2276})} has a left derived bifunctor 
\[ (- \ot^{\mrm{L}}_A -) 
: \dcat{D}(B) \times \dcat{D}(B) \to \dcat{D}(B \ot_{A} B) . \]
If either $M_0$ or $M_1$ is a complex of $B$-modules that is K-flat over $A$, 
then the morphism 
$\eta^{\mrm{L}}_{M_0, M_1} : M_0 \ot^{\mrm{L}}_A M_1 \to M_0 \ot_A M_1$
in $\dcat{D}(B \ot_{A} B)$ is an isomorphism. 
\end{lem}

\begin{proof}
This is a variant of Theorem \ref{thm:2107}. 
We know by Corollary \ref{cor:1580} and Proposition \ref{prop:1525} 
that every complex $M \in \dcat{C}(B)$ admits a K-flat resolution 
$P \to M$. Because $B$ is flat over $A$, the complex $P$ is also K-flat over 
$A$. By Theorem \ref{thm:2106} the left derived bifunctor 
$(- \ot^{\mrm{L}}_A -)$ exists, and the condition on 
$\eta^{\mrm{L}}_{M_0, M_1}$ holds. 
\end{proof}

\begin{rem} \label{rem:2275}
The innocuous looking Lemma \ref{lem:2276} is actually of 
tremendous importance. Without the flatness of $A \to B$ we could do very 
little homological algebra of bimodules. Getting around the lack of flatness 
requires the use of flat DG ring resolutions, as explained in Remark 
\ref{rem:2285}. 
\end{rem}

Every module $L \in \dcat{M}(B)$ has an action by $B \ot_{A} B$ coming from the 
homomorphism $\ep$ in (\ref{eqn:2281}). Consider now a module 
$N \in \dcat{M}(B \ot_{A} B)$.
The $\K$-module $N$ has two possible $B$-module structures, coming from the 
homomorphisms $\eta_i$. Thus the $\K$-module
$\opn{Hom}_{B \ot_{A} B}(L, N)$
has three possible $B$-module structures: there is one action from the 
$B$-module structure on $L$, and there are two from the 
$B$-module structures on $N$. 

\begin{lem} \label{lem:2325}
The three $B$-module structures on $\opn{Hom}_{B \ot_{A} B}(L, N)$ coincide. 
\end{lem}

\begin{exer} \label{exer:2325}
Prove the lemma. 
\end{exer}

We are mostly interested in the $B$-module $L = B$. As the module $N$ changes,
we get an additive functor 
\[ \opn{Hom}_{B \ot_{A} B}(B, -) : \dcat{M}(B \ot_{A} B) \to \dcat{M}(B) . \]
Passing to complexes, and then to homotopy categories, we get a triangulated 
functor 
\[ \opn{Hom}_{B \ot_{A} B}(B, -) : \dcat{K}(B \ot_{A} B) \to \dcat{K}(B) . \]
This has a right derived functor
\begin{equation} \label{eqn:2279}
\opn{RHom}_{B \ot_{A} B}(B, -) :
\dcat{D}(B \ot_{A} B) \to \dcat{D}(B)  , 
\end{equation}
that is calculated by K-injective resolutions. Namely if 
$I \in \dcat{C}(B \ot_{A} B)$ is a K-injective complex, then the morphism 
\[ \eta^{\mrm{R}}_{B, I} : \opn{Hom}_{B \ot_{A} B}(B, I) \to 
\opn{RHom}_{B \ot_{A} B}(B, I) \]
in $\dcat{D}(B)$ is an isomorphism. 

By composing the bifunctor $(- \ot^{\mrm{L}}_A -)$ from Lemma \ref{lem:2276} 
and the functor \lb 
$\opn{RHom}_{B \ot_{A} B}(B, -)$ from (\ref{eqn:2279})
we obtain a triangulated bifunctor 
\begin{equation} \label{eqn:2280}
\opn{RHom}_{B \ot_{A} B}(B, - \ot^{\mrm{L}}_{A} -) :
\dcat{D}(B) \times \dcat{D}(B) \to \dcat{D}(B) .
\end{equation}

\begin{dfn} \label{dfn:2275} 
Under Setup \tup{\ref{set:2285}}, the {\em squaring operation} 
\index{Squaring operation}
\index{1-SqBA@$\opn{Sq}_{B / A}$}
is the functor 
\[ \opn{Sq}_{B / A} : \dcat{D}(B) \to \dcat{D}(B) \]
defined as follows:
\begin{enumerate}
\item For a complex $M \in \dcat{D}(B)$, its square is the complex 
\[ \opn{Sq}_{B / A}(M) := 
\opn{RHom}_{B \ot_{A} B}(B, M \ot^{\mrm{L}}_{A} M) \in \dcat{D}(B) . \]

\item For a morphism 
$\phi : M \to N$ in $\dcat{D}(B)$, its square is the morphism 
\[ \opn{Sq}_{B / A}(\phi) := 
\opn{RHom}_{B \ot_{A} B}(B, \phi \ot^{\mrm{L}}_{A} \phi) :
\opn{Sq}_{B / A}(M) \to \opn{Sq}_{B / A}(N) \]
in $\dcat{D}(B)$.
\end{enumerate}
\end{dfn}

It will be good to have an explicit formulation of the squaring operation. 
Let us first choose a K-projective resolution 
$\si : P \to M$ in $\dcat{C}(B)$. 
Note that $P$ is unique up to homotopy equivalence, and $\si$ is unique up to 
homotopy. Since $B$ is flat over $A$, the complex $P$ is K-flat over $A$. We 
get an isomorphism 
$\opn{pres}_{P} : P \ot_A P \iso M \ot^{\mrm{L}}_{A} M$ 
in $\dcat{D}(B \ot_{A} B)$, that we call a {\em presentation}. It is uniquely 
characterized by the commutativity of the diagram
\[ \UseTips \xymatrix @C=10ex @R=6ex {
M \ot^{\mrm{L}}_{A} M
\ar[r]^{\eta^{\mrm{L}}_{M, M}}
&
M \ot_A M
\\
P \ot^{\mrm{L}}_{A} P
\ar[r]_{\eta^{\mrm{L}}_{P, P}}^{\cong}
\ar[u]^{\opn{Q}(\si) \ot^{\mrm{L}}_A \opn{Q}(\si)}_{\cong}
&
P \ot_A P
\ar[ul]_{\opn{pres}_{P}}^(0.6){\cong}
\ar[u]_{\opn{Q}(\si \ot_A \si)}
} \]
in $\dcat{D}(B \ot_{A} B)$. 

Next we choose a K-injective resolution 
$\rho : P \ot_A P \to I$ in $\dcat{C}(B \ot_{A} B)$. 
The complex $I$ is unique up to homotopy equivalence,
and the homomorphism $\rho$ is unique up to homotopy. The resolution 
$\rho$ gives rise to an isomorphism
\[ \opn{pres}_{I} : \opn{Hom}_{B \ot_{A} B}(B, I) \iso 
\opn{RHom}_{B \ot_{A} B}(B, P \ot_{A} P) \]
in $\dcat{D}(B)$ such that the diagram 
\[ \UseTips \xymatrix @C=10ex @R=7ex {
\opn{Hom}_{B \ot_{A} B}(B, P \ot_{A} P)
\ar[r]^{\eta^{\mrm{R}}_{B, P \ot_{A} P}}
\ar[d]_{\opn{Q}(\opn{Hom}(\opn{id}, \rho))}
&
\opn{RHom}_{B \ot_{A} B}(B, P \ot_{A} P)
\ar[d]^{\opn{RHom}(\opn{id}, \opn{Q}(\rho))}_{\cong}
\\
\opn{Hom}_{B \ot_{A} B}(B, I)
\ar[r]_{\eta^{\mrm{R}}_{B, I}}^{\cong}
\ar[ur]^{\opn{pres}_{I}}_(0.6){\cong}
&
\opn{RHom}_{B \ot_{A} B}(B, I)
} \]
is commutative. 

The combination of the presentations $\opn{pres}_{P}$ and 
$\opn{pres}_{I}$ gives an isomorphism 
\[ \opn{pres}_{P, I} : 
\opn{Hom}_{B \ot_{A} B}(B, I) \iso \opn{Sq}_{B / A}(M) \]
in $\dcat{D}(B)$, that we also call a presentation. 

Let $\phi : M \to N$ be a morphism in $\dcat{D}(B)$. The morphism 
$\opn{Sq}_{B / A}(\phi)$ can also be made explicit using presentations. 
For that we need to choose a K-projective resolution 
$\si_N : Q \to N$ in $\dcat{C}(B)$, and a K-injective resolution 
$\rho_N : Q \ot_A Q \to J$ in $\dcat{C}(B \ot_{A} B)$. 
These provide us with a presentation 
\[ \opn{pres}_{Q, J} : \opn{Hom}_{B \ot_{A} B}(M, J) \iso \opn{Sq}_{B / A}(N) .
\]
There are homomorphisms
$\til{\phi} : P \to Q$ in $\dcat{C}_{\mrm{str}}(B)$,
and 
$\chi : I \to J$ in $\dcat{C}_{\mrm{str}}(B \ot_{A} B)$,
both unique up to homotopy, such that the diagrams 
\[ \UseTips \xymatrix @C=8ex @R=6ex {
P
\ar[r]^{\opn{Q}(\si)}_{\cong}
\ar[d]_{\opn{Q}(\til{\phi})}
&
M
\ar[d]^{\phi}
\\
Q
\ar[r]^{\opn{Q}(\si_N)}_{\cong}
&
N
} 
\qquad \quad 
\UseTips \xymatrix @C=9ex @R=6ex {
M \ot^{\mrm{L}}_A M
\ar[d]_{\phi \ot^{\mrm{L}}_A \phi}
&
P \ot_A P
\ar[l]_{\opn{pres}_{P}}^{\cong}
\ar[r]^(0.6){\opn{Q}(\rho)}_{\cong}
\ar[d]_{\opn{Q}(\til{\phi} \ot_A \til{\phi})}
&
I
\ar[d]^{\opn{Q}(\chi)}
\\
N \ot^{\mrm{L}}_A N
&
Q \ot_A Q
\ar[l]_{\opn{pres}_{Q}}^{\cong}
\ar[r]^(0.6){\opn{Q}(\rho_N)}_(0.6){\cong}
&
J
}
\]
in $\dcat{D}(B)$ and $\dcat{D}(B \ot_{A} B)$ respectively are commutative. 
See Subsections \ref{subsec:K-inj} and \ref{subsec:K-proj}. Then the diagram 
\begin{equation} \label{eqn:2329}
\UseTips \xymatrix @C=10ex @R=6ex {
\opn{Hom}_{B \ot_{A} B}(B, I)
\ar[r]^(0.55){\opn{pres}_{P, I}}_(0.55){\cong}
\ar[d]_{\opn{Q}(\opn{Hom}(\opn{id}, \chi))}
&
\opn{Sq}_{B / A}(M)
\ar[d]^{\opn{Sq}_{B / A}(\phi)}
\\
\opn{Hom}_{B \ot_{A} B}(B, J)
\ar[r]^(0.55){\opn{pres}_{Q, J}}_(0.55){\cong}
&
\opn{Sq}_{B / A}(N)
} 
\end{equation}
in $\dcat{D}(B)$ is commutative.

The squaring operation is not an additive functor. In fact, it is a 
{\em quadratic functor}:

\begin{thm} \label{thm:2325} 
Let $\phi : M \to N$ be a morphism in $\dcat{D}(B)$ and let $b \in B$.
Then 
\[ \opn{Sq}_{B / A}(b \cd \phi) = b^2 \cd \opn{Sq}_{B / A}(\phi) , \]
as morphisms 
$\opn{Sq}_{B / A}(M) \to \opn{Sq}_{B / A}(N)$ in $\dcat{D}(B)$. 
\end{thm}

\begin{proof}
We shall use presentations. Let 
$\til{\phi} : P \to Q$ be a homomorphism in $\dcat{C}_{\mrm{str}}(B)$
that represents $\phi$, as above. Then the homomorphism
$b \cd \til{\phi} : P \to Q$ 
$\dcat{C}_{\mrm{str}}(B)$ represents $b \cd \phi$.
Tensoring we get a homomorphism
\[ (b \cd \til{\phi}) \ot_A (b \cd \til{\phi}) : 
P \ot_A P \to Q \ot_A Q \]
$\dcat{C}_{\mrm{str}}(B \ot_A B)$.
But 
\[ (b \cd \til{\phi}) \ot_A (b \cd \til{\phi}) = 
(b \ot b) \cd (\til{\phi} \ot_A \til{\phi}) . \]
Hence on the K-injectives we get the homomorphism
$(b \ot b) \cd \chi : I \to  J$ 
$\dcat{C}_{\mrm{str}}(B \ot_A B)$.  We conclude that 
\[ \opn{Hom}_{B \ot_{A} B} \bigl( \opn{id}_B, (b \ot b) \cd \chi \bigr) : 
\opn{Hom}_{B \ot_{A} B}(B, I) \to 
\opn{Hom}_{B \ot_{A} B}(B, J) \]
represents 
$\opn{Sq}_{B / A}(b \cd \phi)$. Finally, by Lemma \ref{lem:2325} 
we know that 
\[ \begin{aligned}
& \opn{Hom}_{B \ot_{A} B} \bigl( \opn{id}_B, (b \ot b) \cd \chi \bigr) = 
\opn{Hom}_{B \ot_{A} B}(b^2 \cd \opn{id}_B, \chi) 
\\
& \quad = b^2 \cd  \opn{Hom}_{B \ot_{A} B}(\opn{id}_B, \chi) . 
\end{aligned} \]
\end{proof}

\begin{dfn} \label{dfn:2276} 
Let $M \in \dcat{D}(B)$. 
A {\em rigidifying isomorphism} 
\index{Rigidifying isomorphism}
for $M$ over $B$ 
relative to $A$ is an isomorphism 
$\rho : M \iso \opn{Sq}_{B / A}(M)$
in $\dcat{D}(B)$. 
\end{dfn}

\begin{dfn} \label{dfn:2295}
A {\em rigid complex} 
\index{Rigid complex}
over $B$ relative to $A$ is a pair 
$(M, \rho)$, consisting of a complex $M \in \dcat{D}(B)$ 
and a rigidifying isomorphism 
$\rho : M \iso \opn{Sq}_{B / A}(M)$
in $\dcat{D}(B)$. 
\end{dfn}

\begin{dfn} \label{dfn:2296}
Suppose $(M, \rho)$ and $(N, \si)$ are rigid complexes over $B$ relative 
to $A$. A {\em morphism of rigid complexes}
\index{Rigid complex! morphism of {\indash}s} 
$\phi : (M, \rho) \to (N, \si)$ 
is a morphism $\phi : M \to N$ in $\dcat{D}(B)$, such that the diagram 
\[ \UseTips \xymatrix @C=8ex @R=6ex {
M
\ar[r]^(0.35){\rho}_(0.35){\cong}
\ar[d]_{\phi}
&
\opn{Sq}_{B / A}(M)
\ar[d]^{\opn{Sq}_{B / A}(\phi)}
\\
N
\ar[r]^(0.35){\si}_(0.35){\cong}
&
\opn{Sq}_{B / A}(N)
} \]
in $\dcat{D}(B)$ is commutative. 

The category of rigid complexes over $B$ relative to $A$ is denoted by
\index{1-D(A)rigK@$\dcat{D}(A)_{\mrm{rig} / \K}$}
$\dcat{D}(B)_{\mrm{rig} / A}$. 
\end{dfn}

Recall (Definition \ref{dfn:2394}) that a complex $M \in \dcat{D}(B)$ 
has the derived Morita property 
if the derived homothety morphism 
$\opn{hm}^{\mrm{R}}_{M} : B \to \opn{RHom}_B(M, M)$
in $\dcat{D}(B)$ is an isomorphism. 

\begin{thm} \label{thm:2275}
Let $(M, \rho)$ be a rigid complex over $B$ relative to $A$. If $M$ has the 
derived Morita property, then the only automorphism of 
$(M, \rho)$ in $\dcat{D}(B)_{\mrm{rig} / A}$ is the identity. 
\end{thm}

\begin{proof}
Let 
$\phi : (M, \rho) \iso (M, \rho)$
be an automorphism in $\dcat{D}(B)_{\mrm{rig} / A}$.
By Proposition \ref{prop:2335}, there is a unique invertible element $b \in B$ 
such that $\phi = b \cd \opn{id}_M$, as morphisms $M \to M$ in 
$\dcat{D}(B)$. 

Next, according to Theorem \ref{thm:2325}, we have 
\[ \opn{Sq}_{B / A}(\phi) = \opn{Sq}_{B / A}(b \cd \opn{id}_M) = 
b^2 \cd \opn{Sq}_{B / A}(\opn{id}_M) . \]
Plugging this into the diagram in Definition \ref{dfn:2296} we get a 
commutative diagram 
\[ \UseTips \xymatrix @C=8ex @R=6ex {
M
\ar[r]^(0.35){\rho}_(0.35){\cong}
\ar[d]_{b \cd \opn{id}_M}
&
\opn{Sq}_{B / A}(M)
\ar[d]^{b^2 \cd \opn{id}_M}
\\
M
\ar[r]^(0.35){\rho}_(0.35){\cong}
&
\opn{Sq}_{B / A}(M)
} \]
in $\dcat{D}(B)$. 
Once more using Proposition \ref{prop:2335} we see that $b^2 = b$. 
Because $b$ is an invertible element, it follows that $b = 1$. Thus 
$\phi = \opn{id}_M$. 
\end{proof}

\begin{exa} \label{exa:2275}
Assume $B = A$, and take $M := B$. Then $B \ot_{A} B \cong B$, 
$M \ot^{\mrm{L}}_A M \cong M$, and there are canonical isomorphisms
\[ \opn{Sq}_{B / A}(M) = \opn{RHom}_{B \ot_{A} B}(B, M \ot^{\mrm{L}}_A M) \cong 
\opn{Hom}_{B}(B, M) \cong M . \]
Thus the pair $(M, \opn{id}_M)$ belongs to $\dcat{D}(B)_{\mrm{rig} / A}$.
Furthermore, the complex $M$ has the derived Morita property, so Theorem 
\ref{thm:2275} applies. 
\end{exa}

\begin{rem} \label{rem:3120}
To the reader who might object to Example \ref{exa:2275} for being silly, 
we have two things to say. First, the existence theorem for rigid dualizing 
complexes \ref{thm:3210}, whose proof is beyond the scope of this book, starts  
with the {\em tautological rigid structure} of $\K \in \dcat{D}(\K)$, where 
$\K$ is the regular base ring. The rigid structure is then propagated to all 
flat essentially finite type (FEFT) $\K$-rings using {\em induction} and {\em 
coinduction} of rigid structures. See the exercise and examples below for an 
indication of these procedures. 

The second fact we want to point out is that when $A$ is an FEFT $\K$-ring, 
then {\em $A$ has exactly one rigid complex} 
$(M, \rho) \in \dcat{D}(A)_{\mrm{rig} / \K}$ (up to a unique isomorphism) 
that is nonzero on each connected component of $\opn{Spec}(A)$; and it is the 
rigid dualizing complex of $A$ (that will be defined in the next subsection). 
This is proved in \cite{YeZh3}, \cite{AIL} and \cite{Sha3}. 
\end{rem}

\begin{exer} \label{exer:2395}
Assume $A \neq 0$, and take $B := A[t_1, \ldots, t_n]$, the polynomial ring in 
$n$ variables. 
\begin{enumerate}
\item Let $C := B \ot_A B$, and let $I$ be the kernel of the multiplication 
homomorphism $\ep : C \to B$. Show that $I$ is generated by the sequence 
$\bsym{c} = (c_1, \ldots, c_n)$, where 
$c_j := t_j \ot 1 - 1 \ot t_j$. 

\item Show that the Koszul complex 
$\opn{K}(C; \bsym{c})$ is a free resolution of $B$ over $C$. 
(Hint: for $1 \leq m \leq n$ define 
$B_m := A[t_1, \ldots, t_m]$ and 
$C_m := B_m \ot_A B_m$. 
By  direct calculation show that 
$\opn{K}(C_1; \bsym{c}_1) \to B_1$
is a quasi-isomorphism. Next, there is a ring isomorphism 
$B_m \ot_A B_1 \iso B_{m + 1}$, 
$1 \ot t_1 \mapsto t_{m + 1}$.
This induces a ring isomorphism 
$C_m \ot_A C_1 \cong C_{m + 1}$ 
and an isomorphism of complexes 
\[ \opn{K}(C_m; \bsym{c}_m) \ot_A \opn{K}(C_1; \bsym{c}_1)
\cong \opn{K}(C_{m + 1}; \bsym{c}_{m + 1}) . \]
Use this last formula and induction on $m$ to prove that 
$\opn{K}(C_m; \bsym{c}_m) \to B_m$ 
is a quasi-isomorphism.)

\item  Prove that 
\[ \opn{Ext}^i_{C}(B, C) \cong 
\begin{cases}
B & \tup{if} \quad i = n \\
0 & \tup{if} \quad i \neq n  . 
\end{cases} \]
(Hint: use the Koszul resolution from above.)

\item Conclude from item (3) that the complex $B[n] \in \dcat{D}(B)$ is rigid 
relative to $A$; namely that there is a rigidifying isomorphism 
$\rho : B[n] \iso \opn{Sq}_{B / A} \bigl( B[n] \bigr)$
in $\dcat{D}(B)$. 
\end{enumerate}
\end{exer}

\begin{exa} \label{exa:2576}
Let $A$ be a nonzero noetherian ring, and let
$B := \lb A[t_1, \ldots, t_n]$ as in the exercise above. 
Let $\Delta_{B / A} := \Om^n_{B / A}$, the module of degree $n$ 
differential forms of $B$ relative to $A$. 
It is the $n$-th exterior power of $\Om^1_{B / A}$, so it is a free 
$B$ module of rank $1$, with basis 
$\d(t_1) \wedge \cdots \wedge \d(t_n)$.
By the previous example, there exists some rigidifying isomorphism $\rho$ for 
the complex $\Delta_{B / A}[n]$. 

However, it can be shown (see \cite{YeZh3} or \cite{Ye11}) that the complex
$\Delta_{B / A}[n]$ has a {\em canonical rigidifying isomorphism relative to 
$A$}. I.e.\ there is a rigidifying isomorphism 
$\rho : \Delta_{B / A}[n] \iso 
\opn{Sq}_{B / A} \bigl( \Delta_{B / A}[n] \bigr)$
in $\dcat{D}(B)$ which is invariant under $A$-ring automorphisms of $B$. 

This rigidifying isomorphism $\rho$ is an incarnation of Grothendieck's 
Fundamental Local Isomorphism \cite[Proposition II.7.2]{RD} and the Residue 
Isomorphism
\cite[Theorem II.8.2]{RD}. 
\end{exa}

It is well-known that a finitely generated module $M$ over a noetherian ring 
$A$ is flat iff it is projective. See \cite[Corollary to Theorem 7.12]{Mats}. 

\begin{exa} \label{exa:2575}
Let $A$ be a noetherian ring, and let $A \to B$ be a finite flat ring 
homomorphism. So $B$ is a finitely generated projective $A$-module. Define 
$\Delta_{B / A} := \opn{Hom}_{A}(B, A) \in \dcat{M}(B)$.
It can be shown (see \cite{Ye11}) that the complex 
$\Delta_{B / A}$ has a canonical rigidifying isomorphism relative to $A$. I.e.\ 
there is a rigidifying isomorphism 
$\rho : \Delta_{B / A} \iso \opn{Sq}_{B / A}(\Delta_{B / A})$
in $\dcat{D}(B)$ which is invariant under $A$-ring automorphisms of $B$. 
\end{exa}

\begin{rem} \label{rem:4190}
The squaring operation $\opn{Sq}_{B / A}$ is functorial also in the ring $B$, 
in two directions. For this remark we assume that the ring $A$ is noetherian, 
and the rings $B$ and $C$ are flat essentially finite type $A$-rings. (See 
Remark \ref{rem:2285} regarding the flatness.)

Suppose $u : B \to C$ is a finite 
$A$-ring homomorphism, and $\th : N \to M$ is a backward (or trace) morphism 
in $\dcat{D}(B)$ over $u$, in the sense of Definition \ref{dfn:3110}. 
Then there is a backward morphism
\[ \opn{Sq}_{u / A}(\th) : \opn{Sq}_{C / A}(N) \to \opn{Sq}_{B / A}(M) \]
in $\dcat{D}(B)$ over $u$. The morphism $\opn{Sq}_{u / A}(\th)$ is functorial 
in $u$ and $\th$ in an obvious sense. Furthermore, under suitable finiteness 
conditions, if $\th$ is a nondegenerate backward morphism (see Definition 
\ref{dfn:3190}), then so is $\opn{Sq}_{u / A}(\th)$. We call this property the 
{\em trace functoriality} of the squaring operation. 

The trace functoriality of the squaring operation allows us to talk about {\em 
rigid trace morphisms}. With $u : B \to C$ as above, suppose that 
$(M, \rho) \in \dcat{D}(B)_{\mrm{rig} / A}$ and 
$(N, \si) \in \dcat{D}(C)_{\mrm{rig} / A}$. A {\em rigid trace morphism} 
$\th : (N, \si) \to (M, \rho)$
over $u$ relative to $A$ is a trace morphism $\th : N \to M$ over $u$ such 
that the diagram 
\[ \UseTips \xymatrix @C=8ex @R=6ex {
N
\ar[r]^(0.35){\si}_(0.35){\cong}
\ar[d]_{\th}
&
\opn{Sq}_{C / A}(N)
\ar[d]^{\opn{Sq}_{u / A}(\th)}
\\
M
\ar[r]^(0.35){\rho}_(0.35){\cong}
&
\opn{Sq}_{B / A}(M)
} \]
in $\dcat{D}(B)$ is commutative. 

Next, suppose $v : B \to C$ is an essentially \'etale homomorphism 
(i.e.\ formally \'etale, plus the EFT condition that was already assumed) of 
$A$-rings, and $\la : M \to N$ is a forward morphism over $v$, in 
the sense of Definition \ref{dfn:3111}. Then there is a forward morphism 
\[ \opn{q}_{v / A}(\la) : \opn{Sq}_{B / A}(M) \to \opn{Sq}_{C / A}(N) \]
in $\dcat{D}(B)$ over $u$. The morphism $\opn{Sq}_{v / A}(\la)$ is functorial 
in $v$ and $\la$ in an obvious sense. Furthermore, under suitable finiteness 
conditions, if $\la$ is a nondegenerate forward morphism (see Definition 
\ref{dfn:3191}), then so is $\opn{Sq}_{v / A}(\la)$. We call this property the 
{\em \'etale functoriality} of the squaring operation. 

The \'etale functoriality of the squaring operation allows us to talk about 
{\em rigid localization morphisms}. With $v : B \to C$ as above, suppose that 
$(M, \rho) \in \dcat{D}(B)_{\mrm{rig} / A}$, and 
$(N, \si) \in \dcat{D}(C)_{\mrm{rig} / A}$. A {\em rigid localization morphism} 
$\la : (M, \rho) \to (N, \si)$
over $v$ relative to $A$ is a forward morphism $\la : N \to M$ over $v$ such 
that the diagram 
\[ \UseTips \xymatrix @C=8ex @R=6ex {
M
\ar[r]^(0.35){\rho}_(0.35){\cong}
\ar[d]_{\la}
&
\opn{Sq}_{B / A}(M)
\ar[d]^{\opn{Sq}_{v / A}(\la)}
\\
N
\ar[r]^(0.35){\si}_(0.35){\cong}
&
\opn{Sq}_{C / A}(N)
} \]
in $\dcat{D}(B)$ is commutative. 

These functorialities of the squaring operation are hard to prove. An imprecise 
study of them was made in the papers \cite{ YeZh3} and \cite{YeZh4}. 
A correct treatment of the trace functoriality can be found in the paper 
\cite{Ye9}; and a correct treatment of the \'etale functoriality will 
appear in \cite{Ye11}. 
\end{rem}

\begin{rem} \label{rem:2347}
The squaring operation is related to {\em Hochschild cohomology}.
Assume for simplicity that $A$ is a field and $M$ is a $B$-module. 
Then for each $i$ the cohomology 
\[ \opn{H}^i (\opn{Sq}_{B / A}(M)) = 
\opn{Ext}^i_{B \ot_{A} B}(B, M \ot_A M) \]
is the $i$-th Hochschild cohomology with values in the $B$-bimodule 
$M \ot_A M$. 
For more on this material see \cite{AILN}, \cite{Sha1} and \cite{Sha2}. 
\end{rem}

\begin{rem} \label{rem:2285}
It is possible to avoid the assumption that $B$ is flat over $A$. 
This is done by choosing a commutative DG ring $\til{B}$ that is K-flat as a DG 
$A$-module, with a DG ring quasi-isomorphism $\til{B} \to B$ over $A$. 
Such resolutions always exist (see \cite{YeZh3} or \cite{Ye9}). Then we take 
\begin{equation} \label{eqn:2345}
\opn{Sq}_{B / A}(M) := 
\opn{RHom}_{\til{B} \ot_{A} \til{B}}(B, M \ot^{\mrm{L}}_A M) .
\end{equation}
This was the construction used by Zhang and Yekutieli in the paper 
\cite{YeZh3}.

Unfortunately there was a serious error in \cite{YeZh3}: 
we did not prove that formula (\ref{eqn:2345}) is independent of  
the resolution $\til{B}$. This error was discovered, and corrected,
by L. Avramov, S. Iyengar, J. Lipman and S. Nayak in their paper \cite{AILN}. 

There were ensuing errors in \cite{YeZh3} regarding the functoriality of the 
squaring operation in the ring $B$ (as described in Remark \ref{rem:4190}). 
The paper \cite{AILN} did not treat such functoriality at all, and the 
constructions and proofs of the trace functoriality were corrected only 
in our recent paper \cite{Ye9}. It is worthwhile to mention that the correct 
proofs (both in \cite{AILN} and \cite{Ye9}) rely on {\em noncommutative DG 
rings} and {\em DG bimodules} over them; and the squaring operation is 
replaced by the {\em rectangle operation}. 

Because the non-flat case is so much more complicated, we have decided not to 
reproduce it in the book. The interested reader can look up the research papers 
\cite{Ye9}, \cite{Ye11}, \cite{Ye12} and \cite{Ye13}, the survey article 
\cite{Ye5}, or the lecture notes \cite{Ye10}. 

A general treatment of derived categories of bimodules, based on K-flat DG ring 
resolutions, will be in the paper \cite{Ye14}.
\end{rem}



\mysubsection{Rigid Dualizing and Residue Complexes}
\label{subsec:rig-res-cplx-rng}

In this subsection we make these assumptions:

\begin{setup} \label{set:3210}
$\K$ is a nonzero regular noetherian commutative ring (Definition 
\ref{dfn:3075}). All other rings are in the 
category $\ringcfeft{\K}$ of FEFT commutative $\K$-rings, where 
``FEFT'' is short for ``flat essentially finite type''. 
\end{setup}

\begin{dfn} \label{dfn:2530}
A {\em rigid dualizing complex}% 
\index{Dualizing complex! commutative rigid}
over $A$ relative to $\K$ is a rigid complex 
$(R, \rho)$ over $A$ relative to $\K$, as in Definition \ref{dfn:2295}, such 
that $R$ is a dualizing complex over $A$, in the sense of Definition 
\ref{dfn:2155}.
\end{dfn}

See Remark \ref{rem:2285} regarding the flatness condition. 

\begin{dfn} \label{dfn:4200}
For $A \in \ringcfeft{\K}$ we write  
$A^{\mrm{en}} := A \ot_{\K} A$ for the 
{\em enveloping ring}%
\index{Enveloping ring}%
\index{1-Aen@$A^{\mrm{en}}$}
of $A$ relative to $\K$.
\end{dfn}

With this notation the square of a complex $M \in \dcat{D}(A)$ relative to $\K$ 
is 
\[ \opn{Sq}_{A / \K}(M) = 
\opn{RHom}_{A^{\mrm{en}}}(A, M \otimes^{\mrm{L}}_{\K} M)
\in \dcat{D}(A)  . \]

Recall the category $\dcat{D}(A)_{\mrm{rig} / \K}$ of rigid complexes over 
$A$ relative to $\K$, from Definition \ref{dfn:2296}. 

\begin{thm}[Uniqueness] \label{thm:2280}
Let $A$ be a flat essentially finite type ring over the regular noetherian ring 
$\K$. If $A$ has a rigid dualizing complex $(R, \rho)$, then it is 
unique up to a unique isomorphism in $\dcat{D}(A)_{\mrm{rig} / \K}$. 
\end{thm}

\begin{proof}
Suppose $(R', \rho')$ is another rigid dualizing 
complex over $A$ relative to $\K$. Let 
$A = \prod_{i = 1}^r A_i$ be the connected component decomposition of $A$. 
Corollary \ref{cor:2175} says that 
$R' \cong R \ot^{\mrm{L}}_A P$,
where 
$P \cong \bigoplus_{i = 1}^r L_i[n_i]$
for integers $n_i$ and rank $1$ projective $A_i$-modules $L_i$. 
There is an isomorphism 
\begin{equation} \label{eqn:2545}
R' \otimes^{\mrm{L}}_{\K} R' = 
(R \ot^{\mrm{L}}_A P) \otimes^{\mrm{L}}_{\K} (R \ot^{\mrm{L}}_A P) \cong 
(R \otimes^{\mrm{L}}_{\K} R) \otimes^{\mrm{L}}_{A^{\mrm{en}}}
(P \otimes^{\mrm{L}}_{\K} P) 
\end{equation}
in $\dcat{D}(A^{\mrm{en}})$, and 
$P \otimes^{\mrm{L}}_{\K} P$
has finite flat dimension over $A^{\mrm{en}}$. 
So we have this sequence of isomorphisms in $\dcat{D}(A)$~:
\begin{equation} \label{eqn:2530}
\begin{aligned}
& R \ot^{\mrm{L}}_A P \cong R' \cong \opn{Sq}_{A / \K}(R')
=  \mrm{RHom}_{A^{\mrm{en}}}(A, R' \otimes^{\mrm{L}}_{\K} R') 
\\
& \quad 
\cong^{\diamondsuit} 
\mrm{RHom}_{A^{\mrm{en}}}
\bigl( A, (R \otimes^{\mrm{L}}_{\K} R) \otimes^{\mrm{L}}_{A^{\mrm{en}}}
(P \otimes^{\mrm{L}}_{\K} P) \bigr) 
\\
& \quad 
\cong^{\dag} \mrm{RHom}_{A^{\mrm{en}}}(A, R \otimes^{\mrm{L}}_{\K} R) 
\otimes^{\mrm{L}}_{A^{\mrm{en}}} (P \otimes^{\mrm{L}}_{\K} P)
\\
& \quad 
= \opn{Sq}_{A / \K}(R) 
\otimes^{\mrm{L}}_{A^{\mrm{en}}} (P \otimes^{\mrm{L}}_{\K} P)
\\
& \quad 
\cong R \otimes^{\mrm{L}}_{A^{\mrm{en}}} (P \otimes^{\mrm{L}}_{\K} P)
\cong R \ot^{\mrm{L}}_A P  \ot^{\mrm{L}}_A P .
\end{aligned}
\end{equation}
The isomorphism $\cong^{\diamondsuit}$ is by (\ref{eqn:2545}), and the 
isomorphism $\cong^{\dag}$ is by Theorem \ref{thm:2700}.
We also used the rigidifying isomorphisms of $R$ and $R'$. Now 
\[ \opn{RHom}_{A}(R, R \ot^{\mrm{L}}_A P) \cong 
\opn{RHom}_{A}(R, R) \ot^{\mrm{L}}_A P \cong P , \]
again using Theorem \ref{thm:2700}, and by the derived Morita property of 
$R$. Likewise there are isomorphisms
\[ \begin{aligned}
& 
\opn{RHom}_{A}(R, R \ot^{\mrm{L}}_A P \ot^{\mrm{L}}_A P) \cong 
\opn{RHom}_{A}(R, R \ot^{\mrm{L}}_A P) \ot^{\mrm{L}}_A P
\\
& \quad
\cong \opn{RHom}_{A}(R, R) \ot^{\mrm{L}}_A P \ot^{\mrm{L}}_A P
\cong P \ot^{\mrm{L}}_A P 
\end{aligned} \]
in $\dcat{D}(A)$.
Thus, together with (\ref{eqn:2530}), we deduce that 
$P \ot^{\mrm{L}}_A P \cong P$. 
But then on each connected component $A_i$ we have 
\[ L_i[n_i] \cong L_i[n_i] \ot_A L_i[n_i]  = 
(L_i \ot_A L_i)[2 \cd n_i] . \]
This implies that $L_i \cong A_i$ and $n_i = 0$. 
We see that actually $P \cong A$, so there is an isomorphism 
$\phi_0 : R \iso R'$ in $\dcat{D}(A)$. 

The isomorphism $\phi_0$ might not be rigid; but due to the derived 
Morita property of the complex $R'$, there is an invertible element $a \in A$ 
such that 
$\opn{Sq}_{A / \K}(\phi_0) \circ \rho = a \cd \rho' \circ \phi_0$ 
as isomorphisms 
$R \iso \opn{Sq}_{A / \K}(R')$. 
Define $\phi := a^{-1} \cd \phi_0$. 
Then, according to Theorem \ref{thm:2325}, we have 
\[ \opn{Sq}_{A / \K}(\phi) \circ \rho = 
a^{-2} \cd \opn{Sq}_{A / \K}(\phi_0) \circ \rho = 
a^{-2} \cd a \cd \rho' \circ \phi_0 = \rho' \circ \phi . \]
We see that
$\phi : (R, \rho) \iso (R', \rho')$
is a rigid isomorphism. Its uniqueness is already known by Theorem 
\ref{thm:2275}, since $R'$ has the derived Morita property.
\end{proof}

\begin{thm}[Existence] \label{thm:3210}
Let $A$ be a flat essentially finite type ring over the regular noetherian ring 
$\K$. Then $A$ has a rigid dualizing complex $(R_A, \rho_A)$. 
\end{thm}

The proof of this theorem is beyond the scope of the book, as it requires 
the study of {\em induced and coinduced rigid structures}. 
For a proof see \cite{Ye11}. (There is an incomplete proof in 
\cite[Theorem 1.1]{YeZh4}.)
The next two examples provide proofs in some cases. 

\begin{exa} \label{exa:4200}
Suppose $A = \K[t_1, \ldots, t_n]$, the polynomial ring in $n$ variables.
(More generally, we can take $A$ to be any essentially smooth $\K$-ring of 
relative dimension $n$). Define the complex 
$R_A := \Om^n_{A / \K}[n]$. According to Example \ref{exa:2576} there is a 
rigidifying isomorphism 
$\rho_A : R_A \iso \opn{Sq}_{A / \K}(R_A)$.
Because $A$ is a regular ring, it follows that $R_A$ is a dualizing complex. We 
see that $(R_A, \rho_A)$ is a rigid dualizing complex. 
\end{exa}

\begin{exa} \label{exa:4201}
Suppose $A$ is a finite type $\K$-ring. As explained in Example \ref{exa:4506},
$A$ has a noncommutative rigid dualizing complex, that is actually $A$-central. 
Thus it is a rigid dualizing complex over $A$ relative to $\K$ in the sense of 
Definition \ref{dfn:2530}. 
\end{exa}

The dimension function $\opn{dim}_R$ relative to a dualizing complex $R$
was introduced in Definition \ref{dfn:2200}. If $R' \cong R$, then of course 
the dimension functions satisfy 
$\opn{dim}_{R'} = \opn{dim}_R$. 
In view of the previous theorem, the next definition is valid. 

\begin{dfn} \label{dfn:2537}
Let $A \in \ringcfeft{\K}$. The {\em rigid dimension function 
relative to $\K$} is the function
\[   \opn{rig{.}dim}_{\K} : \opn{Spec}(A) \to \Z \]
given by the formula 
$\opn{rig{.}dim}_{\K}(\p) := \opn{dim}_R(\p)$,
where $R$ is any rigid dualizing complex over $A$ relative to $\K$. 
\end{dfn}

The next examples (taken from \cite{Ye11}) give an idea what the rigid 
dimension function looks like.  

\begin{exa} \label{exa:3235}
Let $\K$ be a field and $A$ a finite type $\K$-ring. 
Then for every $\p \in \opn{Spec}(A)$ there is equality
$\opn{rig{.}dim}_{\K}(\p) = \opn{dim}(A / \p)$, 
where the latter is the Krull dimension of the ring $A / \p$. 
\end{exa}

\begin{exa} \label{exa:3236}
Again $\K$ is a field, but now $L$ is a finitely generated field extension of 
$\K$; or in other words, $L$ is a field in $\ringcfeft{\K}$ . 
Then, writing $\p := (0) \sub L$, we have  
$\opn{rig{.}dim}_{\K}(\p) = \opn{tr{.}deg}_{\K}(L)$,
the transcendence degree of the field extension $\K \to L$. 
\end{exa}

\begin{exa} \label{exa:3237}
Take $\K = A = \Z$. For a maximal ideal $\m = (p) \sub \Z$ we 
have $\opn{rig{.}dim}_{\K}(\m) = -1$; and for the generic ideal 
$\p = (0) \sub \Z$ we have $\opn{rig{.}dim}_{\K}(\p) = 0$. 
\end{exa}

Residue complexes were introduced in Subsection \ref{subsec:res-cplx-ring}.

\begin{dfn} \label{dfn:2535}
A {\em rigid residue complex} 
\index{Residue complex! rigid}
over $A$ relative to $\K$ is a 
rigid complex $(\KK_A, \rho_A)$ over $A$ relative to $\K$, such that $\KK_A$ is 
a residue complex over $A$. 
\end{dfn}

Using the rigid dimension function relative to $\K$, we have this decomposition 
of the $A$-module $\KK^{-i}_A$ for each $i$~:
$\KK_A^{-i} \cong \bigoplus_{\opn{rig{.}dim}_{\K}(\p) = i} J(\p)$,
where $J(\p)$ is the indecomposable injective module corresponding to the prime 
ideal $\p$. 

In Definition \ref{dfn:2296} we introduced the category 
$\dcat{D}(A)_{\mrm{rig} / \K}$. Recall that the objects of 
$\dcat{D}(A)_{\mrm{rig} / \K}$
are rigid complexes $(M, \rho)$ over $A$ relative to $\K$; and the morphisms 
$\phi : (M, \rho) \to (N, \si)$
are the morphisms $\phi : M \to N$ in $\dcat{D}(A)$ for which there is equality 
$\si \circ \phi = \opn{Sq}_{A / \K}(\phi) \circ \rho$.
Rigid residue complexes live, or rather move, in another category. 

\begin{dfn} \label{dfn:2536}
The category $\dcat{C}(A)_{\mrm{rig} / \K}$
\index{1-C(A)rigK@$\dcat{C}(A)_{\mrm{rig} / \K}$}
is defined as follows. Its objects are the rigid complexes $(M, \rho)$ over $A$ 
relative to $\K$. Given two objects 
$(M, \rho)$ and $(N, \si)$, a morphism 
$\phi : (M, \rho) \to (N, \si)$
in $\dcat{C}(A)_{\mrm{rig} / \K}$ is a morphism 
$\phi : M \to  N$ in $\dcat{C}_{\mrm{str}}(A)$,
such that the diagram 
\[ \UseTips \xymatrix @C=8ex @R=6ex {
M
\ar[d]_{\opn{Q}(\phi)}
\ar[r]^(0.35){\rho}_(0.35){\cong}
&
\opn{Sq}_{A / \K}(M)
\ar[d]^{\opn{Sq}_{A / \K}(\opn{Q}(\phi))}
\\
N
\ar[r]^(0.35){\si}_(0.35){\cong}
&
\opn{Sq}_{A / \K}(N)
} \]
in $\dcat{D}(A)$ is commutative.
\end{dfn}

Let us emphasize the hybrid nature of the category 
$\dcat{C}(A)_{\mrm{rig} / \K}$~:
the morphisms are homomorphisms of complexes (literally degree $0$ 
homomorphisms of graded $A$-modules  $\phi : M \to N$ that commute with the 
differentials); but they must satisfy a compatibility condition (rigidity) in 
the derived category. 

\begin{thm} \label{thm:2535}
Let $A$ be an FEFT ring over the regular noetherian ring 
$\K$. The ring $A$ has a rigid residue complex $(\KK_{A}, \rho_{A})$ 
relative to $\K$, and it is unique, up to a unique isomorphism in 
$\dcat{C}(A)_{\mrm{rig} / \K}$.
\end{thm}

\begin{proof}
Existence: by Theorem \ref{thm:3210} there is a rigid dualizing complex 
$(R_{A}, \rho_{A})$ over $A$ relative to $\K$. Let 
$\KK_{A}$ be the minimal injective resolution of the complex $R_{A}$. 
According to Theorem \ref{thm:2245}, $\KK_{A}$ is a residue complex. It 
inherits the rigidifying isomorphism $\rho_{A}$ from $R_{A}$. So the 
pair $(\KK_{A}, \rho_{A})$ is a rigid residue complex over $A$ relative to 
$\K$. 

Uniqueness: suppose $(\KK', \rho')$ is another rigid residue complex over 
$A$ relative to $\K$.  Theorem \ref{thm:2280} tells us that there is a unique 
isomorphism 
$\phi : (\KK_{A}, \rho_{A}) \iso (\KK', \rho')$
in $\dcat{D}(A)_{\mrm{rig} / \K}$. 
Now the dimension functions of these two residue complexes are equal (both are 
$\opn{rig{.}dim}_{\K}$). So by Theorem \ref{thm:2240} the function 
\[ \opn{Q} : \opn{Hom}_{\dcat{C}(A)_{\mrm{rig} / \K}}
\bigl( (\KK_{A}, \rho_{A}), (\KK', \rho') \bigr) \to 
\opn{Hom}_{\dcat{D}(A)_{\mrm{rig} / \K}}
\bigl( (\KK_{A}, \rho_{A}), (\KK', \rho') \bigr) \]
is bijective. Thus $\phi$ lifts uniquely to an isomorphism in 
$\dcat{C}(A)_{\mrm{rig} / \K}$. 
\end{proof}

\begin{rem} \label{rem:4191}   
The rigid residue complex $\KK_{A}$ is functorial in the ring $A$ in two 
different ways, that we briefly explain here. As in Setup \ref{set:3210}, we 
work in the category $\ringcfeft{\K}$. The flatness condition can be removed 
(at a price, see Remark \ref{rem:2285}), but it seems that the assumption that 
the base ring $\K$ is regular, and the other rings are EFT over it, are 
necessary.  
 
Suppose $u : A \to B$ is a finite homomorphism in $\ringcfeft{\K}$. 
Then there is a homomorphism 
$\opn{tr}_{B / A} : \KK_B \to \KK_A$
in $\dcat{C}_{\mrm{str}}(A)$, called the {\em rigid trace homomorphism}. 
It is a nondegenerate rigid trace morphism over $u$ (see Remark \ref{rem:4190}),
namely it respects the rigidifying isomorphisms, and the induced homomorphism 
$\KK_B \to \opn{Hom}_{A}(B, \KK_A)$
in $\dcat{C}_{\mrm{str}}(B)$ is an isomorphism. Furthermore, generalizing 
Theorem \ref{thm:2535}, $\opn{tr}_{B / A}$ is the unique 
nondegenerate rigid trace morphism over $u$. If $B \to C$ is another finite 
homomorphism in $\ringcfeft{\K}$, then 
$\opn{tr}_{B / A} \circ \opn{tr}_{C / B}  = \opn{tr}_{C / A}$.

Suppose $v : A \to A'$ is an essentially \'etale homomorphism in 
$\ringcfeft{\K}$. Then there is a homomorphism 
$\opn{q}_{A' / A} : \KK_{A} \to \KK_{A'}$
in $\dcat{C}_{\mrm{str}}(A)$, called the {\em rigid localization homomorphism}. 
It is a nondegenerate rigid localization morphism over $v$ (see Remark 
\ref{rem:4190}), namely it respects the rigidifying isomorphisms, and the 
induced homomorphism 
$A' \ot_{A} \KK_{A} \to \KK_{A'}$
in $\dcat{C}_{\mrm{str}}(A')$ is an isomorphism. Furthermore,
$\opn{q}_{A' / A}$ is the unique nondegenerate rigid localization morphism over 
$v$. If $A' \to A''$ is another essentially \'etale homomorphism in 
$\ringcfeft{\K}$, then 
$\opn{q}_{A'' / A'} \circ \opn{q}_{A' / A} = \opn{q}_{A'' / A}$.

The two functorialities of the rigid residue complex $\KK_A$ commute with each 
other, in the following sense. Suppose we are given the first commutative 
diagram below in $\ringcfeft{\K}$, which is cocartesian (i.e.\ 
$B' \cong A' \ot_A B$), $A \to B$ is finite, and $A \to A'$ is essentially 
\'etale. Then the second diagram in $\dcat{C}_{\mrm{str}}(A)$ is commutative. 
\[ 
\UseTips \xymatrix @C=6ex @R=6ex {
A
\ar[r]^{}
\ar[d]_{}
&
B
\ar[d]
\\
A' 
\ar[r]
&
B'
}
\qquad \quad 
\UseTips \xymatrix @C=8ex @R=6ex {
\KK_{A}
\ar[d]_{\opn{q}_{A' / A}}
&
\KK_{B}
\ar[l]_{\opn{tr}_{B / A}}
\ar[d]^{\opn{q}_{B' / B}}
\\
\KK_{A'}
&
\KK_{B'}
\ar[l]_{\opn{tr}_{B' / A'}}
} 
\]

Imprecise proofs of these assertions can be found in the papers \cite{ YeZh3} 
and \cite{YeZh4}. A correct treatment will appear in \cite{Ye11}. 
\end{rem}

\begin{rem} \label{rem:4192}  
The functorialities of the rigid residue complexes that were outlined in Remark 
\ref{rem:4190} enable similar construction for schemes. In this remark we shall 
dispense with the flatness condition (invoking Remark \ref{rem:2285}); so  
$\K$ is a nonzero regular noetherian base ring (e.g.\ a field or the ring of 
integers $\Z$), and we consider the category $\catt{Sch} \eftover \K$
of EFT $\K$-schemes. 

Let $X$ be an EFT $\K$-scheme. A {\em rigid residue complex} on $X$ is a pair
$(\KK_X, \bsym{\rho}_X)$, consisting of a bounded complex of injective 
quasi-coherent $\OO_X$-modules $\KK_X$, and a {\em rigid structure} 
$\bsym{\rho}_X$ on $\KK_X$, to be explained below. 
For every affine open set $U \sub X$ which is 
{\em strictly EFT}, namely $A := \Ga(U, \OO_X)$ is an EFT $\K$-ring, the 
complex $\KK_A := \Ga(U, \KK_X)$
is equipped with a rigidifying isomorphism $\rho_U$, making it into a 
rigid residue complex over $A$ relative to $\K$.
If $U' \sub U$ is a smaller strictly EFT open set of $X$, 
then the restriction homomorphism 
$\KK_{A} \to \KK_{A'}$
must be the unique rigid localization homomorphism, for the given 
rigidifying isomorphisms $\rho_{U}$ and $\rho_{U'}$.
The collection of these rigidifying isomorphisms 
$\bsym{\rho}_X := \{ \rho_U \}$
is the rigid structure on $\KK_X$. 
The scheme $X$ admits a rigid residue complex $(\KK_X, \bsym{\rho}_X)$, and it 
is unique up to a unique rigid isomorphism. 

If $f : Y \to X$ is a finite map in $\catt{Sch} \eftover \K$, then 
there is the rigid trace homomorphism 
$\opn{tr}_f : f_{*}(\KK_Y) \to \KK_X$
in $\dcat{C}_{\mrm{str}}(X) := \dcat{C}_{\mrm{str}}(\cat{Mod} \OO_X)$; 
and it is nondegenerate, in the sense that the 
induced homomorphism 
$f_{*}(\KK_Y) \to \Hom_{X} \bigl( f_*(\OO_Y), \KK_X \bigr)$
in $\dcat{C}_{\mrm{str}}(X)$ is an isomorphism. 
If $g : X' \to X$ is an essentially \'etale map in 
$\catt{Sch} \eftover \K$, then 
there is the rigid localization homomorphism 
$\opn{q}_g : \KK_X \to g_{*}(\KK_{X'})$
in $\dcat{C}_{\mrm{str}}(X)$; and it is nondegenerate, in the sense that the 
induced homomorphism 
$g^*(\KK_X) \to \KK_{X'}$
is an isomorphism in $\dcat{C}_{\mrm{str}}(X')$. In this way the rigid residue 
complexes become a sheaf on the small \'etale site of $X$. 

Actually, an infinitesimal local construction gives, for every morphism 
$f : Y \to X$ $\catt{Sch} \eftover \K$, 
the {\em ind-rigid trace homomorphism}
$\opn{tr}_f : f_{*}(\KK_Y) \to \KK_X$
in $\dcat{G}_{\mrm{str}}(X)$; namely this is a homomorphism of graded 
quasi-coherent sheaves on $X$. The {\em Residue Theorem} says that when $f$ is 
proper, the homomorphism $\opn{tr}_f$ commutes with the differentials. 
And the {\em Duality Theorem} says that in the proper case, the ind-rigid trace 
$\opn{tr}_f$ induces an isomorphism 
\[ \mrm{R} f_{*} \bigl( \mrm{R} \Hom_{Y}(\NN, \KK_Y) \bigr) \to  
\mrm{R} \Hom_{X} \bigr( \mrm{R} f_{*}(\NN), \KK_X \bigr) \]
in $\dcat{D}(X)$ for every 
$\NN \in \dcat{D}_{\mrm{c}}^{\mrm{b}}(X)$.
This explicit rigid version of the original theorems of Grothendieck from 
\cite{RD} will appear in \cite{Ye12}. 
\end{rem}


\begin{rem} \label{rem:4193}
In this remark we consider {\em Deligne-Mumford stacks}, of finite type 
over the regular noetherian base ring $\K$. Due to the \'etale functoriality of 
the rigid residue complexes, as explained in Remark \ref{rem:4191}, every such 
stack $\mfrak{X}$ admits a rigid residue complex 
$(\KK_{\mfrak{X}}, \bsym{\rho}_{\mfrak{X}})$.
Here the rigid structure 
$\bsym{\rho}_{\mfrak{X}} = \{ \rho_{(U, g)} \}$
is indexed by the \'etale maps $g : U \to \mfrak{X}$ from strictly EFT affine 
$\K$-schemes $U$. 

To a map $f : \mfrak{Y} \to \mfrak{X}$ between DM stacks 
we associate the ind-rigid trace 
$\opn{tr}_{f} : f_{*}(\KK_{\mfrak{Y}}) \to \KK_{\mfrak{X}}$,
which is a homomorphism of graded quasi-coherent sheaves on $\mfrak{X}$. 
In order to construct the ind-rigid trace for stacks we need another property 
of rigid residue complexes over rings, that we call {\em \'etale codescent}. 
Suppose $v : A \to A'$ is a faithfully \'etale ring homomorphism of EFT 
$\K$-rings. Let $w_1, w_2 : A' \to A' \ot_{A} A'$ be the two inclusions. 
\'Etale codescent says that in every degree $i$ the sequence of $A$-module 
homomorphisms 
\[ \KK^i_{A' \ot_{A} A'} \xar{ \, \opn{tr}_{w_1} - \opn{tr}_{w_2} \, }
\KK^i_{A'} \xar{ \, \opn{tr}_{v} \, } \KK^i_A \to 0 \]
is exact. 

As for the Residue Theorem: we can only prove it for a proper map of DM stacks 
$f : \mfrak{Y} \to \mfrak{X}$ which is {\em coarsely schematic}.
Likewise, we can only prove the Duality Theorem for a proper map of DM stacks 
$f$ which is coarsely schematic and {\em tame}. These conditions, and a sketch 
of the proofs, can be found in the lecture notes \cite{Ye10}. 
These results shall be written in detail in the future paper \cite{Ye13}. 
Observe that duality is not expected to hold without the tameness condition; 
but the coarsely schematic condition seems to be a temporary technical hitch. 
\end{rem}
  







